\clearpage
\onecolumn
\thispagestyle{empty}
\vspace*{\fill}
\begin{center}
	{\color{radblue}\rule{\linewidth}{2pt}}\par
	\vspace{6mm}
	{\Huge\bfseries\color{radblue}Columna}\par
	\vspace{6mm}
	{\color{radblue}\rule{\linewidth}{2pt}}
\end{center}
\vspace*{\fill}
\clearpage
\twocolumn
\pagestyle{contenido}



%%%%%%%%%%%%%%%%%%%%%%%%%%%%%%%%%%%%%%%%%%%%%%%%%%%%%%%%%%%
%  SECCIÓN: Columna Cervical — Articulación Atlanto-Axoidea
%  Archivo: secciones/columna_cervical.tex
%  Sin preámbulo ni \begin{document} — lo maneja main.tex
	%%%%%%%%%%%%%%%%%%%%%%%%%%%%%%%%%%%%%%%%%%%%%%%%%%%%%%%%%%%
	
	\section{Columna Cervical}
	
	\subsection{Articulación Atlanto-Axoidea - Enfoque Transbucal (Boca Abierta)}
	
	\subsubsection*{1. Identificación}
	
	\textbf{Tipo:} Proyección complementaria.
	
	\textbf{Indicaciones clínicas principales:}
	\begin{itemize}
		\item Traumatismo cervical alto: fractura de Jefferson (C1) y fractura de apófisis odontoides (C2)
		\item Fracturas de masas laterales de C1
		\item Cervicalgia, artritis reumatoidea con sospecha de subluxación del ligamento transverso
		\item Pacientes politraumatizados por accidentes de tránsito
		\item Luxaciones atlanto-axoideas e inestabilidad C1-C2
	\end{itemize}
	
	\subsubsection*{2. Técnica Base}
	
	\textbf{Receptor de imagen:} 18$\times$24 cm o 24$\times$30 cm. Orientación \textbf{longitudinal}, centrado a la altura del axis.
	
	\textbf{DFRI:} 70 cm (preferencial). Se aprovecha la divergencia del rayo: con el tubo más cerca, el haz permite proyectar la base del cráneo hacia arriba y los dientes fuera del campo, despejando la odontoides en el centro del receptor. Puede utilizarse hasta 1 metro, pero siempre menos de ese valor.
	
	\textbf{Técnica:} Intermedia (homogeneizada). La boca actúa como ventana natural, lo que resulta suficiente para obtener detalle de la odontoides, valorar el trabeculado óseo y detectar líneas de fractura cuando existieran.
	
	\textbf{Bucky:} Bucky de mesa con grilla.
	
	\subsubsection*{3. Posicionamiento}
	
	\paragraph{Paciente}
	Decúbito supino. Brazos a los lados del cuerpo; puede colocarse un apoyo bajo las rodillas para mayor comodidad. El plano medio sagital debe estar perpendicular y centrado en línea media respecto al receptor.
	
	\begin{figure}[h]
		\centering
		\includegraphics[width=0.7\linewidth]{imagenes/col1}
		\caption{}
		\label{fig:col1}
	\end{figure}
	
	
	\paragraph{Región de interés}
	Se ajusta la cabeza de manera que el \textbf{plano de oclusión} —línea trazada entre el borde inferior de los incisivos superiores y el vértice de la apófisis mastoides— quede \textbf{perpendicular al receptor de imagen}. Los hombros deben permanecer en el mismo plano horizontal. Se le pide al paciente que abra la boca al máximo.
	
	\paragraph{Rayo Central}
	Perpendicular al plano del receptor. Punto de entrada: \textbf{punto medio de la boca abierta}.
	
	\paragraph{Receptor de imagen}
	Centrado a la altura del axis (C2). La región a incluir abarca desde el borde inferior de los incisivos superiores hasta el borde superior del mentón (para asegurar el cuerpo de C2), y lateralmente las ramas mandibulares. Colimación ajustada a los cuatro lados del área de interés, aproximadamente 10$\times$10 cm.
	
	\paragraph{Respiración}
	Durante la exposición, indicar al paciente que mantenga la boca abierta y diga suavemente \textit{``ah''} de forma sostenida. La fonación hace que la lengua descienda y quede quieta, evitando su superposición sobre C1-C2 y previniendo la generación de imágenes falsas (pseudofracturas). Se le indica además que no respire ni trague durante la exposición.
	
	\subsubsection*{4. Anatomía Visible}
	
	Se visualizan la apófisis odontoides (C2), el atlas (C1) con sus masas laterales, el cuerpo del axis y las articulaciones atlanto-axoideas (C1-C2) a través de la cavidad bucal. Deben verse todas las superficies articulares del atlas y el axis para permitir la comprobación de desplazamiento lateral. Las ramas mandibulares deben proyectarse equidistantes respecto a la odontoides. La sombra de la lengua no debe superponerse sobre el atlas ni el axis.
	
	\subsubsection*{5. Criterios de Evaluación}
	
	\textbf{Posicionamiento correcto:} Odontoides visible sin superposición en el centro de la imagen. Borde inferior de los incisivos superiores y la base del cráneo alineados y ubicados por encima de la punta de la odontoides. Ramas mandibulares equidistantes. Boca bien abierta. Sombra de la lengua no proyectada sobre C1-C2. La posición es óptima cuando las sombras de la superficie de oclusión de los incisivos centrales superiores, la base del cráneo y los vértices de las apófisis mastoides se encuentran en línea.
	
	\textbf{Ausencia de rotación:} El espacio entre la odontoides y las masas laterales de C1 debe ser igual de ambos lados. Si un lado se observa más ancho que el otro, el paciente está rotado. El plano medio sagital (línea que pasa por nariz y mentón) debe quedar perpendicular al receptor sin desviación lateral.
	
	\textbf{Penetración y contraste:} Técnica intermedia que permite valorar el trabeculado óseo de la odontoides y detectar líneas de fractura finas. La imagen debe mostrar contraste suficiente para diferenciar las corticales de C1 y C2 del espacio articular y los tejidos blandos circundantes.
	
	\subsubsection*{6. Puntos Críticos}
	
	\textbf{Incisivos superiores superpuestos a la odontoides (mentón muy bajo):} Los dientes tapan la odontoides y pueden simular una línea de fractura (pseudofractura). Corrección: subir levemente el mentón aumentando la extensión de la cabeza.
	
	\textbf{Base del cráneo (occipital) superpuesta a la odontoides (mentón muy elevado):} La nuca cubre la odontoides impidiendo su valoración. Corrección: bajar levemente el mentón reduciendo la extensión.
	
	\textbf{Adaptaciones:} En pacientes con traumatismo cervical no debe intentarse ningún movimiento sin autorización médica previa. Antes de realizar este estudio debe haberse obtenido una radiografía lateral con rayo central horizontal para descartar inestabilidad. Si no se logra un enfoque adecuado luego de dos intentos, se debe recurrir a métodos alternativos como el enfoque de Fuchs o de Judd.
	
	\textbf{Trucos:} La fonación sostenida con \textit{``ah''} es el recurso más efectivo para inmovilizar la lengua y evitar pseudofracturas. La DFRI de 70 cm (menor de 1 metro) es clave: aprovecha la divergencia del haz para despejar la base del cráneo y los dientes del campo de la odontoides. Recordar siempre que este es un \textbf{enfoque complementario} que requiere correlación con la proyección lateral cervical.
	
	\onecolumn
	\clearpage
	
	\begin{figure}
		\centering
		\includegraphics[width=0.7\linewidth]{imagenes/col2}
		\caption{Fractura en masa lateral}
		\label{fig:col2}
	\end{figure}
	
	\twocolumn
	\clearpage
	
	%----------------------------------------------------------
	%  Para agregar otro enfoque duplicá el bloque
	%  desde \subsection hasta \textbf{Trucos}
	%----------------------------------------------------------
	
	
	%%%%%%%%%%%%%%%%%%%%%%%%%%%%%%%%%%%%%%%%%%%%%%%%%%%%%%%%%%%
	%  SECCIÓN: Columna Cervical
	%  Archivo: secciones/columna_cervical.tex
	%  Sin preámbulo ni \begin{document} — lo maneja main.tex
		%%%%%%%%%%%%%%%%%%%%%%%%%%%%%%%%%%%%%%%%%%%%%%%%%%%%%%%%%%%
		
	\section{Columna Cervical}
		
		\subsection{Odontoides AP - Método de Fuchs}
		
		\subsubsection*{1. Identificación}
		
		\textbf{Tipo:} Proyección complementaria.
		
		\textbf{Indicaciones clínicas principales:}
		\begin{itemize}
			\item Artrosis severa que impida la apertura bucal para la proyección transoral.
			\item Trismo (limitación de la apertura bucal).
			\item Paciente con collarín cervical que no puede abrir la boca.
			\item Demostración de la odontoides cuando la mitad superior no se visualiza con claridad en boca abierta.
			\item Inestabilidad atlanto-axoidea.
		\end{itemize}
		
		\textbf{Contraindicación absoluta:} No intentar esta posición ante sospecha de trauma agudo. La extensión forzada del cuello puede desplazar una fractura inestable y causar daño neurológico permanente. Las fracturas de odontoides (Jefferson) están contraindicadas.
		
		\subsubsection*{2. Técnica Base}
		
		\textbf{Receptor de imagen:} 18$\times$24 cm (o 24$\times$30 cm). Orientación \textbf{transversal} centrado a la altura de los vértices de las apófisis mastoides.
		
		\textbf{DFRI:} 100 cm. Distancia estándar para proyecciones de columna cervical.
		
		\textbf{Técnica:} Homogeneizada. El rayo central debe atravesar la base del cráneo, estructura de alta densidad; se requiere un kVp elevado para penetrar y homogeneizar las densidades y visualizar la odontoides por transparencia con nitidez.
		
		\textbf{Bucky:} Bucky de mesa con grilla.
		
		\subsubsection*{3. Posicionamiento}
		
		\paragraph{Paciente}
		Decúbito supino sobre la mesa. Brazos extendidos a los lados del cuerpo. Hombros ajustados para que queden en el mismo plano horizontal. Se puede colocar un apoyo bajo las rodillas para mayor comodidad del paciente.
		
		\begin{figure}[h]
			\centering
			\includegraphics[width=0.7\linewidth]{imagenes/col3}
			\caption{}
			\label{fig:col3}
		\end{figure}
				
		\paragraph{Región de interés}
		El paciente eleva el mentón hasta que la línea mento-meatal (trazada desde el conducto auditivo externo hasta el mentón) quede perpendicular al receptor de imagen. El plano medio sagital debe quedar perpendicular y centrado respecto al RI.
		
		\paragraph{Rayo Central}
		Perpendicular al plano del receptor. Punto de entrada: \textbf{plano medio sagital, inmediatamente distal (inferior) al vértice del mentón}. Si el paciente no puede lograr que la línea mento-meatal quede perpendicular al RI, se angula el RC en dirección craneal de modo que siga incidiendo por debajo del vértice del mentón, compensando la falta de extensión.
		
		\paragraph{Receptor de imagen}
		Centrado a la altura de los vértices de las apófisis mastoides. La colimación se ajusta por los cuatro lados a la región C1--C2. La región a incluir comprende: borde superior del agujero occipital, cuerpo completo de C2 y ángulos mandibulares (para verificar el centrado).
		
		\paragraph{Respiración}
		Suspendida en espiración, sin deglutir. Se instruye al paciente a no respirar ni tragar durante la exposición para evitar cualquier movimiento que genere borrosidad en la imagen.
		
		\subsubsection*{4. Anatomía Visible}
		
		La estructura principal a demostrar es la \textbf{apófisis odontoides de C2 completa}, enmarcada dentro del agujero occipital (agujero magno), que actúa como \textbf{ventana radiológica}: espacio delimitado que forma una cavidad a través de la cual se valora la odontoides por transparencia. Se visualizan también la relación espacial entre la odontoides y el agujero magno, y los ángulos mandibulares como referencias de simetría y centrado.
		
		\subsubsection*{5. Criterios de Evaluación}
		
		\textbf{Posicionamiento correcto:} Toda la apófisis odontoides debe estar enmarcada dentro del agujero magno circular. El RC ha incidido inmediatamente distal al vértice del mentón en el PMS.
		
		\textbf{Ausencia de rotación:} Los ángulos mandibulares deben aparecer equidistantes respecto a la apófisis odontoides y respecto a la línea media del RI. El lado que se vea más ancho corresponde al lado menos apoyado. La nariz y el mentón deben estar alineados con la línea media del chasis.
		
		\textbf{Penetración y contraste:} Técnica homogeneizada que permita ver la odontoides con nitidez por detrás de la base del cráneo; densidades equilibradas entre las estructuras óseas superpuestas.
		
		\subsubsection*{6. Puntos Críticos}
		
		\textbf{Falta de extensión (mentón cubre la odontoides):} El mentón superpone la odontoides impidiendo su visualización. Corrección: indicar al paciente que eleve más el mentón, o angular el RC en dirección craneal para que incida igualmente distal al vértice del mentón.
		
		\textbf{Exceso de extensión (base del cráneo cubre la odontoides):} La base del cráneo se superpone sobre la odontoides. Corrección: indicar al paciente que baje el mentón hasta recuperar la posición correcta de la línea mento-meatal.
		
		\textbf{Adaptaciones:} En pacientes con limitación para extender el cuello (collarín, rigidez, posquirúrgico), compensar angulando el RC cranealmente en lugar de forzar la extensión cervical. No intentar la proyección en trauma agudo con sospecha de fractura inestable.
		
		\textbf{Trucos:} La ventana radiológica del agujero magno es el concepto clave de esta proyección: no se visualiza la odontoides directamente, sino \textit{por transparencia} a través del agujero occipital. Recordar que la línea mento-meatal (CAE--mentón) perpendicular al RI es el eje de posicionamiento central, y que la angulación del RC es el recurso corrector cuando el paciente no puede lograrlo anatómicamente.
		
		
		\onecolumn
		\clearpage
		
		\begin{figure}
			\centering
			\includegraphics[width=0.7\linewidth]{imagenes/col4}
			\caption{}
			\label{fig:col4}
		\end{figure}
		
		
		\twocolumn
		\clearpage
		%----------------------------------------------------------
		%  Para agregar otro enfoque duplicá el bloque
		%  desde \subsection hasta \textbf{Trucos}
		%----------------------------------------------------------
	
	%%%%%%%%%%%%%%%%%%%%%%%%%%%%%%%%%%%%%%%%%%%%%%%%%%%%%%%%%%%
	%  SECCIÓN: Columna Cervical
	%  Archivo: secciones/columna_cervical.tex
	%  Sin preámbulo ni \begin{document} — lo maneja main.tex
		%%%%%%%%%%%%%%%%%%%%%%%%%%%%%%%%%%%%%%%%%%%%%%%%%%%%%%%%%%%
		
	\section{Columna Cervical}
		
		\subsection{Odontoides PA - Método de Judd}
		
		\subsubsection*{1. Identificación}
		
		\textbf{Tipo:} Proyección complementaria. Posición inversa al Método de Fuchs; ambos métodos son alternativas cuando el paciente no puede realizar la proyección transoral.
		
		\textbf{Indicaciones clínicas principales:}
		\begin{itemize}
			\item Demostración y evaluación de la apófisis odontoides
			\item Fracturas de odontoides y estructuras óseas circundantes de C1
			\item Alternativa al Método de Fuchs en pacientes con limitaciones (ej.\ trauma, collarín cervical)
			\item Alternativa a la proyección transoral cuando el paciente no puede abrir la boca
		\end{itemize}
		
		\subsubsection*{2. Técnica Base}
		
		\textbf{Receptor de imagen:} 18$\times$24 cm (o 24$\times$30 cm). Orientación \textbf{transversal} para incluir la región C1--C2 en su totalidad.
		
		\textbf{DFRI:} 100 cm. Distancia estándar para proyecciones de columna cervical.
		
		\textbf{Bucky:} Bucky de mesa. Técnica homogeneizada.
		
		\subsubsection*{3. Posicionamiento}
		
		\paragraph{Paciente}
		Decúbito prono sobre la mesa. Manos apoyadas sobre la mesa a los lados de la cabeza. El plano medio sagital debe quedar perpendicular al receptor para evitar rotación.
		
		\begin{figure}[h]
			\centering
			\includegraphics[width=0.7\linewidth]{imagenes/col5}
			\caption{}
			\label{fig:col5}
		\end{figure}
		
		
		\paragraph{Región de interés}
		El mentón descansa directamente sobre la superficie de la mesa. Se extiende la cabeza hasta que la línea mento-meatal (LMM) quede casi perpendicular a la mesa. La LMM une el mentón con el meato auditivo externo. El rayo central puede ajustarse según las necesidades del paciente para mantenerse paralelo a la LMM. Verificar ausencia de rotación de cabeza y cuello.
		
		\paragraph{Rayo Central}
		Paralelo a la línea mento-meatal (LMM). Punto de entrada: \textbf{a través del hueso occipital medio, aproximadamente 2,5 cm inferoposterior a las puntas de las apófisis mastoides y ángulos mandibulares}. Angulación referencial de 20° en vista lateral según necesidad del paciente.
		
		\paragraph{Receptor de imagen}
		Centrado en la región C1--C2. Colimación ajustada por los cuatro lados para incluir toda la odontoides enmarcada en el agujero magno.
		
		\paragraph{Respiración}
		Suspendida. Instruir al paciente que no respire ni trague durante la exposición, para evitar artefactos de movimiento en la región cervical superior.
		
		\subsubsection*{4. Anatomía Visible}
		
		Apófisis odontoides (proceso odontoideo), atlas (C1) y axis (C2). Relación atlanto-axoidea. La odontoides debe aparecer enmarcada dentro del agujero magno. Estructuras de referencia visibles: mandíbula con sus ángulos (gonios) y la V mandibular, cráneo y vértebras cervicales superiores.
		
		\subsubsection*{5. Criterios de Evaluación}
		
		\textbf{Posicionamiento correcto:} La apófisis odontoides queda completamente enmarcada dentro del agujero magno, visible en su totalidad sin superposiciones óseas que la oculten.
		
		\textbf{Ausencia de rotación:} Simetría bilateral de la mandíbula, del cráneo y de las vértebras cervicales; los ángulos mandibulares (gonios) deben verse equidistantes respecto al plano medio. La asimetría indica rotación de cabeza o cuello.
		
		\textbf{Penetración y contraste:} Técnica homogeneizada adecuada para visualizar la odontoides diferenciada del occipucio y las estructuras circundantes.
		
		\subsubsection*{6. Puntos Críticos}
		
		\textbf{Rotación de la cabeza o el cuello:} Genera asimetría mandibular y vertebral; la odontoides se desplaza del centro del agujero magno. Corregir reposicionando al paciente con el plano medio sagital estrictamente perpendicular al RI.
		
		\textbf{LMM no perpendicular a la mesa:} Si la extensión es insuficiente, la mandíbula se superpone sobre la odontoides. Ajustar el rayo central para que quede paralelo a la LMM o aumentar la extensión cervical dentro de las posibilidades del paciente.
		
		\textbf{Adaptaciones:} En pacientes con trauma o collarín cervical que no pueden abrir la boca (proyección transoral contraindicada), el Método de Judd en decúbito prono y el Método de Fuchs en decúbito supino son las alternativas de elección. Puede realizarse también en bipedestación.
		
		\textbf{Trucos:} La diferencia fundamental con el Método de Fuchs es la posición del paciente: Judd en decúbito prono (PA) y Fuchs en decúbito supino (AP). En el Judd el mentón contacta directamente la mesa. El método transoral sigue siendo el ideal, ya que permite ver además las masas laterales del atlas; Judd y Fuchs se usan cuando la boca cerrada es la única opción disponible.
		
		
			\onecolumn
		\clearpage
		
		\begin{figure}
			\centering
			\includegraphics[width=0.7\linewidth]{imagenes/col4}
			\caption{}
			\label{fig:col4}
		\end{figure}
		
		
		\twocolumn
		\clearpage
		
		%----------------------------------------------------------
		%  Para agregar otro enfoque duplicá el bloque
		%  desde \subsection hasta \textbf{Trucos}
		%----------------------------------------------------------
		
		
	%%%%%%%%%%%%%%%%%%%%%%%%%%%%%%%%%%%%%%%%%%%%%%%%%%%%%%%%%%%
	%  SECCIÓN: Columna Cervical
	%  Archivo: secciones/columna_cervical.tex
	%  Sin preámbulo ni \begin{document} — lo maneja main.tex
		%%%%%%%%%%%%%%%%%%%%%%%%%%%%%%%%%%%%%%%%%%%%%%%%%%%%%%%%%%%
		
	\section{Columna Cervical}
		
		\subsection{Frente de Columna Cervical (AP)}
		
		\subsubsection*{1. Identificación}
		
		\textbf{Tipo:} Proyección básica. Forma par con el perfil de columna cervical (RC horizontal).
		
		\textbf{Indicaciones clínicas principales:}
		\begin{itemize}
			\item Cervicalgia y rectificación cervical
			\item Hernia del núcleo pulposo (HNP) con radiculopatía cervical
			\item Traumatismo de tránsito (latigazo cervical, fractura de apófisis espinosa — ``cavador de arcilla'')
			\item Fracturas por compresión vertebral
		\end{itemize}
		
		\subsubsection*{2. Técnica Base}
		
		\textbf{Receptor de imagen:} 18$\times$24 cm (o 24$\times$30 cm). Orientación \textbf{longitudinal} adaptada según biotipo del paciente.
		
		\textbf{DFRI:} 100 cm. Distancia estándar para proyecciones axiales de columna, reduce distorsión geométrica y mantiene definición.
		
		\textbf{Técnica:} Homogeneizada (escala larga). Se emplea suficiente kVp para visualizar simultáneamente tejidos blandos (vía aérea, laringe), detalle óseo (trabeculado y cortical) y espacios intervertebrales en una sola imagen.
		
		\textbf{Bucky:} Sí. Se utiliza Bucky mural (bipedestación) o Bucky de mesa (decúbito supino).
		
		\subsubsection*{3. Posicionamiento}
		
		\paragraph{Paciente}
		Decúbito supino o bipedestación (descalzo). Espalda apoyada contra el receptor. Hombros ubicados en el mismo plano horizontal y paralelos al RI para evitar curvatura lateral aparente de la columna. El plano medio sagital debe ser perpendicular al receptor.
		
		\begin{figure}[h]
			\centering
			\includegraphics[width=0.7\linewidth]{imagenes/col7}
			\caption{}
			\label{fig:col7}
		\end{figure}
				
		\paragraph{Región de interés}
		Se eleva ligeramente el mentón para superponer el occipital con el mentón y evitar que la mandíbula se proyecte sobre las vértebras cervicales. Diferencia de posición occipital-mentón: aproximadamente 0,5 cm. En pacientes con lordosis pronunciada se provee apoyo bajo la cabeza. Referencia práctica: 2 dedos sobre el occipital y 1 dedo por delante del mentón.
		
		\paragraph{Rayo Central}
		Angulado \textbf{15--20° craneal} al plano del receptor. La angulación permite que el haz pase tangencial a los espacios intervertebrales, compensando la lordosis cervical. Punto de entrada: \textbf{C4 (por encima de la nuez de Adán)}. Si se logra correcta superposición anatómica occipital-mentón se angula 15°; si no se logra, se angula 20° (mayor angulación desplaza estructuras y puede comprometer la superposición ya lograda).
		
		\paragraph{Receptor de imagen}
		Centrado en C4. El borde superior se ubica a nivel del conducto auditivo externo (CAE) y el borde inferior debe incluir hasta T2--T3 (pasaje torácico). La región abarca de C3 a T2--T3.
		
		\paragraph{Respiración}
		Suspendida. Se indica al paciente que no respire y no degluta durante la exposición, para evitar movimiento a nivel de la vía aérea (laringe, tráquea) que genere borrosidad en las estructuras blandas del cuello.
		
		\subsubsection*{4. Anatomía Visible}
		
		Se visualizan los cuerpos vertebrales de C3 a C7 y T1--T3 con sus espacios discales intervertebrales abiertos. El Atlas (C1) y la mayor parte del Axis (C2) quedan superpuestos por el occipital y la mandíbula. Las apófisis transversas y articulares aparecen superpuestas entre sí. Son visibles la laringe, la tráquea y tejidos blandos prelaríngeos gracias a la técnica homogeneizada. Las clavículas deben verse paralelas entre sí y simétricas respecto a la columna.
		
		\subsubsection*{5. Criterios de Evaluación}
		
		\textbf{Posicionamiento correcto:} Espacios intervertebrales abiertos y nítidos gracias a la angulación craneal. Mandíbula y base del cráneo superpuestos a C1--C2. Cuerpos vertebrales de C3 hacia abajo claramente individualizados. Clavículas paralelas y simétricas respecto a la columna.
		
		\textbf{Ausencia de rotación:} Apófisis espinosas equidistantes y alineadas con la línea media de los cuerpos vertebrales. Ángulos mandibulares y apófisis mastoides equidistantes entre sí. Cualquier desviación de las apófisis espinosas indica rotación de cabeza o tronco; corregir con el PMS.
		
		\textbf{Penetración y contraste:} Técnica homogeneizada: escala larga con kVp suficiente para demostrar simultáneamente tejido blando (vía aérea, laringe) y detalle óseo (trabeculado, cortical, espacios intervertebrales) en una sola imagen.
		
		\subsubsection*{6. Puntos Críticos}
		
		\textbf{Superposición mandibular excesiva:} El mentón muy bajo proyecta los gonios mandibulares sobre las vértebras superiores, ocultando C3. Se identifica porque la mandíbula aparece cuadrada con gonios visibles. Corrección: elevar el mentón hasta lograr superposición occipital-mentón adecuada.
		
		\textbf{Occipital demasiado posterior:} El hueso occipital se proyecta redondeado sobre la columna (imagen en V), tapando C3. Corrección: ajustar posición del occipital y mentón en diferencia de 0,5 cm.
		
		\textbf{Adaptaciones:} En trauma de tránsito el estudio se realiza en decúbito supino sin movilizar al paciente. El complementario transoral es prioritario cuando el paciente porta collarín cervical. No se realizan estudios funcionales en contexto de traumatismo agudo.
		
		\textbf{Trucos:} La angulación cefálica dirige el haz entre cuerpos vertebrales superpuestos por la lordosis. Contar vértebras desde C7 hacia arriba (la prominente es el punto de referencia), ya que C1--C2 quedan bajo el cráneo. Si la colimación llega a T3--T4 puede reducirse. Para rectificación cervical o cervicalgia el estudio se complementa con funcionales (hiperextensión e hiperflexión) como proyecciones complementarias.
		
		\onecolumn
		\clearpage
		
		\begin{figure}
			\centering
			\includegraphics[width=0.7\linewidth]{imagenes/col6}
			\caption{}
			\label{fig:col6}
		\end{figure}
		
		
		\twocolumn
		\clearpage
		
		
		
		%----------------------------------------------------------
		%  Para agregar otro enfoque duplicá el bloque
		%  desde \subsection hasta \textbf{Trucos}
		%----------------------------------------------------------
		
		
	%%%%%%%%%%%%%%%%%%%%%%%%%%%%%%%%%%%%%%%%%%%%%%%%%%%%%%%%%%%
	%  SECCIÓN: Columna Cervical
	%  Archivo: secciones/columna_cervical.tex
	%  Sin preámbulo ni \begin{document} — lo maneja main.tex
		%%%%%%%%%%%%%%%%%%%%%%%%%%%%%%%%%%%%%%%%%%%%%%%%%%%%%%%%%%%
		
	\section{Columna Cervical}
		
		\subsection{Proyección Lateral de Columna Cervical (Método de Grandy)}
		
		\subsubsection*{1. Identificación}
		
		\textbf{Tipo:} Proyección básica. Forma par con la proyección anteroposterior de columna cervical.
		
		\textbf{Indicaciones clínicas principales:}
		\begin{itemize}
			\item Cervicalgia y discopatías (sospecha de hernia de disco / hernias del núcleo pulposo con pinzamiento radicular)
			\item Rectificación cervical y espondiloartrosis (visualización de osteofitos y desgaste articular)
			\item Traumatismo de tránsito: latigazo, fractura de apófisis espinosa (fractura del cavador de arcilla)
			\item Parestesias (hormigueos) de origen cervical
		\end{itemize}
		
		\subsubsection*{2. Técnica Base}
		
		\textbf{Receptor de imagen:} 18$\times$24 cm (puede utilizarse 24$\times$30 cm). Orientación \textbf{longitudinal} (vertical). El borde superior del receptor se ubica 2,5 cm por encima del conducto auditivo externo (C.A.E.).
		
		\textbf{DFRI:} 140 cm (1,52--1,83 m). La distancia aumentada compensa la separación entre la columna cervical y el receptor impuesta por la presencia de los hombros: los rayos llegan más paralelos y menos divergentes, se reduce la magnificación y se minimiza la distorsión geométrica. El incremento de distancia requiere multiplicar los mAs por 2 respecto a la técnica estándar.
		
		\textbf{Técnica:} Homogeneizada (escala larga). Suficiente kVp para demostrar simultáneamente el tejido blando (vías aéreas, laringe) y el detalle óseo (trabeculado, cortical) junto con los espacios intervertebrales en una sola imagen.
		
		\textbf{Bucky:} Bucky con grilla.
		
		\subsubsection*{3. Posicionamiento}
		
		\paragraph{Paciente}
		Bipedestación o sentado en posición lateral estricta. El plano medio sagital queda paralelo al receptor. Los hombros se desplazan hacia abajo tanto como sea posible (en la práctica: pedir al paciente que tome sus propias manos por detrás de la espalda y baje los hombros, o utilizar pesos en las muñecas). Ambos hombros deben quedar en el mismo plano horizontal para evitar rotación del tronco.
		
		\begin{figure}[h]
			\centering
			\includegraphics[width=0.7\linewidth]{imagenes/col8}
			\caption{}
			\label{fig:col8}
		\end{figure}
		
		
		\paragraph{Región de interés}
		La barbilla se eleva para que las ramas y los ángulos de la mandíbula no se superpongan a C1--C2. El paciente fija la mirada en un punto de la pared para mantener la posición y evitar rotación cervical. El plano medio sagital permanece paralelo al receptor durante toda la exposición.
		
		\paragraph{Rayo Central}
		Horizontal y perpendicular al plano del receptor. Punto de entrada: \textbf{C4, a nivel de la parte superior del cartílago tiroides}.
		
		\paragraph{Receptor de imagen}
		Centrado en C4. El borde superior queda 2,5 cm por encima del C.A.E. La colimación abarca 4 cm a cada lado del cuello. La región incluida comprende los cuerpos vertebrales desde C1 hasta C7, con al menos un tercio de T1 visible.
		
		\paragraph{Respiración}
		Suspendida al final de una espiración máxima. Este momento logra el máximo descenso de los hombros, favoreciendo la visualización de C7 y T1 sin superposición.
		
		\subsubsection*{4. Anatomía Visible}
		
		Las siete vértebras cervicales en perfil estricto y al menos un tercio de T1. Cuerpos vertebrales con morfología lateral sin doble contorno. Apófisis espinosas de perfil. Pilares articulares. Espacios discales intervertebrales abiertos y radiolúcidos. Cinco articulaciones interapofisarias inferiores abiertas y superpuestas entre sí (indicativo de ausencia de rotación). Tejidos blandos adyacentes: vías aéreas, laringe.
		
		\subsubsection*{5. Criterios de Evaluación}
		
		\textbf{Posicionamiento correcto:} Barbilla elevada de modo que las ramas y los ángulos mandibulares no se superpongan a C1--C2. C4 en el centro de la imagen. Las ramas mandibulares se proyectan superpuestas o casi superpuestas entre sí (un leve desfase es normal por diferencia de magnificación entre la rama más cercana y la más lejana al receptor). Apófisis espinosas de perfil.
		
		\textbf{Ausencia de rotación:} Articulaciones interapofisarias superpuestas y abiertas. Cuerpos vertebrales de perfil sin doble contorno. Si existe rotación, las ramas mandibulares aparecen desplazadas: una posterior y otra anterior. Si existe inclinación lateral, una rama queda superior y la otra inferior.
		
		\textbf{Penetración y contraste:} Escala larga (alto kVp) que permite visualizar simultáneamente el detalle óseo (trabeculado y cortical) y los tejidos blandos cervicales (vías aéreas, laringe) en una sola imagen. Espacios discales intervertebrales perceptibles y radiolúcidos.
		
		\subsubsection*{6. Puntos Críticos}
		
		\textbf{Elevación de hombros:} Impide la visualización de C7 y T1 por superposición. Corrección: pedir al paciente que baje los hombros activamente, utilizar pesos en las muñecas o indicar que tome sus manos por detrás; si C7 no es visible, obtener una proyección cervicotorácica separada.
		
		\textbf{Rotación cervical:} Las ramas mandibulares se proyectan desfasadas (una anterior y otra posterior) y los cuerpos vertebrales muestran doble contorno. Corrección: verificar que el plano medio sagital sea estrictamente paralelo al receptor y pedir al paciente que mire un punto fijo en la pared.
		
		\textbf{Adaptaciones:} Si el paciente no puede bajar los hombros voluntariamente, se utilizan pesos sujetos a las muñecas. En pacientes con inclinación lateral que cierra los espacios articulares, corregir dejando el plano medio sagital paralelo al receptor (ajustar el PMS).
		
		\textbf{Trucos:} Las ramas mandibulares siempre quedan levemente desfasadas entre sí por diferencia de magnificación (distinta distancia al receptor); este desfase leve es normal y no debe confundirse con rotación. Para contar las vértebras se procede de arriba hacia abajo. El paciente que fija la mirada en un punto fijo de la pared mantiene mejor la posición y reduce la necesidad de repetir la exposición.
		
		\onecolumn
		\clearpage
		
		\begin{figure}
			\centering
			\includegraphics[width=0.7\linewidth]{imagenes/col9}
			\caption{}
			\label{fig:col9}
		\end{figure}
		
		
		\twocolumn
		\clearpage
		
		
		%----------------------------------------------------------
		%  Para agregar otro enfoque duplicá el bloque
		%  desde \subsection hasta \textbf{Trucos}
		%----------------------------------------------------------
		
		
		
%%%%%%%%%%%%%%%%%%%%%%%%%%%%%%%%%%%%%%%%%%%%%%%%%%%%%%%%%%%
%  SECCIÓN: Columna Cervical
%  Archivo: secciones/funcionales_cervical.tex
%  Sin preámbulo ni \begin{document} — lo maneja main.tex
	%%%%%%%%%%%%%%%%%%%%%%%%%%%%%%%%%%%%%%%%%%%%%%%%%%%%%%%%%%%
	
	\section{Columna Cervical}
	
	\subsection{Proyección Funcional - Hiperflexión}
	
	\subsubsection*{1. Identificación}
	
	\textbf{Tipo:} Proyección complementaria. Forma par con la proyección funcional de hiperextensión; ambas deben realizarse juntas como par radiológico. Si el paciente no cuenta con una de las dos, se efectúa el par completo.
	
	\textbf{Indicaciones clínicas principales:}
	\begin{itemize}
		\item Espondilolistesis (desplazamiento de una vértebra sobre otra).
		\item Latigazo cervical agudo o crónico por accidente de tránsito (evaluación de limitación del movimiento por lesión mecánica).
		\item Evaluación preoperatoria por artritis cervical.
		\item Seguimiento de artrodesis (placa, tornillos).
		\item Demostración del movimiento anteroposterior normal o de la ausencia de movimiento por traumatismo o patología; inestabilidad cervical y subluxaciones ocultas.
	\end{itemize}
	
	\subsubsection*{2. Técnica Base}
	
	\textbf{Receptor de imagen:} 24$\times$30 cm (puede usarse 18$\times$24 cm). Orientación \textbf{longitudinal}; en la práctica, a veces transversal para no cortar la apófisis espinosa de C7.
	
	\textbf{DFRI:} 140--150 cm (bibliografía: 152--183 cm). La distancia aumentada compensa la magnificación generada por la mayor distancia objeto--receptor debida al hombro que se interpone entre la columna y el RI.  Una distancia mayor también contribuye a visualizar C7.
	
	\textbf{Técnica:} Homogenizada, escala larga (suficiente kVp). Permite visualizar simultáneamente el tejido blando (vías aéreas, laringe), el detalle óseo (trabeculado y cortical) y los espacios intervertebrales en una sola imagen. \textbf{kVp --- mAs}
	
	\textbf{Bucky:} Bucky de pared con grilla.
	
	\subsubsection*{3. Posicionamiento}
	
	\paragraph{Paciente}
	De pie o sentado en posición lateral verdadera. El plano medio sagital (PMS) paralelo al RI. Los hombros en un mismo plano, apoyados sobre el RI.
	
	\begin{figure}[h]
		\centering
		\includegraphics[width=0.7\linewidth]{imagenes/col10}
		\caption{}
		\label{fig:col10}
	\end{figure}
	
	
	\paragraph{Región de interés}
	El paciente se acerca pegado al RI para que el hombro se apoye en él. Para la hiperflexión se le indica que lleve el mentón hacia el pecho, intentando tocarse el esternón con la barbilla, manteniendo hombros y espalda quietos. No se permite que dé un paso hacia adelante. Esto sitúa las vértebras cervicales en flexión forzada para la primera exposición.
	
	\paragraph{Rayo Central}
	Horizontal, perpendicular al plano del receptor, sin angulación. Punto de entrada: \textbf{C4 (nuez de Adán)}.
	
	\paragraph{Receptor de imagen}
	Centrado a la altura de C4. El borde superior del RI se ubica aproximadamente 5 cm por encima del conducto auditivo externo (CAE). Se colima a los cuatro lados del área de interés. Deben incluirse las vértebras cervicales superiores e inferiores (C1--C7), con especial atención a C7 en hiperflexión, así como las articulaciones interapofisarias y los discos intervertebrales.
	
	\paragraph{Respiración}
	Suspendida en espiración completa: el paciente expulsa todo el aire y mantiene la apnea durante la exposición. El descenso natural de los hombros en espiración favorece la inclusión de C7.
	
	\subsubsection*{4. Anatomía Visible}
	
	Las siete vértebras cervicales (C1--C7) en posición lateral verdadera, con los cuerpos vertebrales y las apófisis espinosas de perfil. Discos intervertebrales y espacios interapofisarios. Articulaciones interapofisarias (carillas articulares) visibles como una única línea continua. Tejidos blandos anteriores: vías aéreas y laringe. En hiperflexión: los cuerpos vertebrales se juntan y las apófisis espinosas se separan (mayor apertura del espacio interespinoso). El cuerpo de la mandíbula adopta una posición casi vertical.
	
	\subsubsection*{5. Criterios de Evaluación}
	
	\textbf{Posicionamiento correcto:} Cuerpo de la mandíbula casi vertical. Las siete apófisis espinosas de perfil. Las siete vértebras cervicales en posición lateral verdadera. Se valora el grado de separación de las apófisis espinosas adyacentes: en hiperflexión normal, dichas apófisis se separan y se alejan de los cuerpos vertebrales.
	
	\textbf{Ausencia de rotación:} Las ramas mandibulares deben aparecer superpuestas entre sí. Si están desfasadas, el paciente presenta rotación: se corrige colocando el PMS paralelo al RI y asegurando que los hombros queden en un mismo plano apoyados sobre el RI. Las articulaciones interapofisarias deben verse como una única línea; si se ven dobles, indica rotación del cuello.
	
	\textbf{Penetración y contraste:} Técnica de escala larga con suficiente kVp que permita ver simultáneamente los tejidos blandos (vías aéreas, laringe) y el detalle óseo (trabeculado, cortical) junto con los espacios intervertebrales en una sola imagen.
	
	\subsubsection*{6. Puntos Críticos}
	
	\textbf{Rotación del paciente:} Las ramas mandibulares aparecen desfasadas y las carillas articulares se ven dobles. Corrección: colocar el PMS estrictamente paralelo al RI y los hombros en un mismo plano apoyados sobre el receptor.
	
	\textbf{Exclusión de C7:} La vértebra C7 queda fuera del campo por descenso insuficiente del hombro o centrado inadecuado. Corrección: pedir espiración completa para descenso natural de hombros; aumentar la DFRI; verificar centrado y campo colimado.
	
	\textbf{Adaptaciones:} No realizar proyecciones funcionales en pacientes con fractura cervical conocida o sospechada. En pacientes con movilidad reducida, registrar el máximo rango alcanzable sin forzar el movimiento.
	
	\textbf{Trucos:} Indicar claramente al paciente que no dé un paso hacia adelante al flexionar, ya que el movimiento del tronco invalida la posición lateral verdadera. La apnea en espiración, además de evitar borrosidad cinética, baja naturalmente los hombros y mejora la visualización de C7.
	
	\onecolumn
	\clearpage
	
	\begin{figure}
		\centering
		\includegraphics[width=0.7\linewidth]{imagenes/col12}
		\caption{}
		\label{fig:col12}
	\end{figure}
	
	
	\twocolumn
	\clearpage
	
	%----------------------------------------------------------
	%  Para agregar otro enfoque duplicá el bloque
	%  desde \subsection hasta \textbf{Trucos}
	%----------------------------------------------------------
	
	\subsection{Proyección Funcional - Hiperextensión}
	
	\subsubsection*{1. Identificación}
	
	\textbf{Tipo:} Proyección complementaria. Forma par con la proyección funcional de hiperflexión; ambas deben realizarse juntas como par radiológico. Si el paciente no cuenta con una de las dos, se efectúa el par completo.
	
	\textbf{Indicaciones clínicas principales:}
	\begin{itemize}
		\item Espondilolistesis (desplazamiento de una vértebra sobre otra).
		\item Latigazo cervical agudo o crónico por accidente de tránsito (evaluación de limitación del movimiento por lesión mecánica).
		\item Evaluación preoperatoria por artritis cervical.
		\item Seguimiento de artrodesis (placa, tornillos).
		\item Demostración del movimiento anteroposterior normal o de la ausencia de movimiento por traumatismo o patología; inestabilidad cervical y subluxaciones ocultas.
	\end{itemize}
	
	\subsubsection*{2. Técnica Base}
	
	\textbf{Receptor de imagen:} 24$\times$30 cm (puede usarse 18$\times$24 cm). Orientación \textbf{longitudinal}; en la práctica, a veces transversal para no cortar la apófisis espinosa de C7.
	
	\textbf{DFRI:} 140--150 cm (bibliografía: 152--183 cm). La distancia aumentada compensa la magnificación generada por la mayor distancia objeto--receptor debida al hombro que se interpone entre la columna y el RI. Una distancia mayor también contribuye a visualizar C7.
	
	\textbf{Técnica:} Homogenizada, escala larga (suficiente kVp). Permite visualizar simultáneamente el tejido blando (vías aéreas, laringe), el detalle óseo (trabeculado y cortical) y los espacios intervertebrales en una sola imagen. \textbf{kVp --- mAs}
	
	\textbf{Bucky:} Bucky de pared con grilla.
	
	\subsubsection*{3. Posicionamiento}
	
	\paragraph{Paciente}
	De pie o sentado en posición lateral verdadera. El plano medio sagital (PMS) paralelo al RI. Los hombros en un mismo plano, apoyados sobre el RI.
	
	\begin{figure}[h]
		\centering
		\includegraphics[width=0.7\linewidth]{imagenes/col11}
		\caption{}
		\label{fig:col11}
	\end{figure}
	
	
	\paragraph{Región de interés}
	El paciente se acerca pegado al RI para que el hombro se apoye en él. Para la hiperextensión se le indica que eleve la barbilla y deje caer la cabeza hacia atrás lo máximo que pueda (como si quisiera mirar el techo detrás suyo), manteniendo hombros y espalda quietos. No se permite que dé un paso hacia atrás. Esto sitúa las vértebras cervicales en extensión forzada para la segunda exposición.
	
	\paragraph{Rayo Central}
	Horizontal, perpendicular al plano del receptor, sin angulación. Punto de entrada: \textbf{C4 (nuez de Adán)}.
	
	\paragraph{Receptor de imagen}
	Centrado a la altura de C4. El borde superior del RI se ubica aproximadamente 5 cm por encima del conducto auditivo externo (CAE). Se colima a los cuatro lados del área de interés. Deben incluirse C1--C7, las articulaciones interapofisarias y los discos intervertebrales.
	
	\paragraph{Respiración}
	Suspendida en espiración completa: el paciente expulsa todo el aire y mantiene la apnea durante la exposición. El descenso natural de los hombros en espiración favorece la inclusión de C7.
	
	\subsubsection*{4. Anatomía Visible}
	
	Las siete vértebras cervicales (C1--C7) en posición lateral verdadera, con los cuerpos vertebrales y las apófisis espinosas de perfil. Discos intervertebrales y espacios interapofisarios. Articulaciones interapofisarias (carillas articulares) visibles como una única línea continua. Tejidos blandos anteriores: vías aéreas y laringe. En hiperextensión: los cuerpos vertebrales se separan y las apófisis espinosas se juntan (mayor cierre del espacio interespinoso). El cuerpo de la mandíbula adopta una posición casi horizontal.
	
	\subsubsection*{5. Criterios de Evaluación}
	
	\textbf{Posicionamiento correcto:} Cuerpo de la mandíbula casi horizontal. Las siete vértebras cervicales en posición lateral verdadera. Se valora el grado de separación de los cuerpos vertebrales adyacentes: en hiperextensión normal, los cuerpos se separan mientras las apófisis espinosas se aproximan.
	
	\textbf{Ausencia de rotación:} Las ramas mandibulares deben aparecer superpuestas entre sí. Si están desfasadas, el paciente presenta rotación: se corrige colocando el PMS paralelo al RI y asegurando que los hombros queden en un mismo plano apoyados sobre el RI. Las articulaciones interapofisarias deben verse como una única línea; si se ven dobles, indica rotación del cuello.
	
	\textbf{Penetración y contraste:} Técnica de escala larga con suficiente kVp que permita ver simultáneamente los tejidos blandos (vías aéreas, laringe) y el detalle óseo (trabeculado, cortical) junto con los espacios intervertebrales en una sola imagen.
	
	\subsubsection*{6. Puntos Críticos}
	
	\textbf{Rotación del paciente:} Las ramas mandibulares aparecen desfasadas y las carillas articulares se ven dobles. Corrección: colocar el PMS estrictamente paralelo al RI y los hombros en un mismo plano apoyados sobre el receptor.
	
	\textbf{Exclusión de C7:} La vértebra C7 queda fuera del campo por descenso insuficiente del hombro o centrado inadecuado. Corrección: pedir espiración completa para descenso natural de hombros; aumentar la DFRI; verificar centrado y campo colimado.
	
	\textbf{Adaptaciones:} No realizar proyecciones funcionales en pacientes con fractura cervical conocida o sospechada. En pacientes con movilidad reducida, registrar el máximo rango alcanzable sin forzar el movimiento.
	
	\textbf{Trucos:} Indicar claramente al paciente que no dé un paso hacia atrás al extender, ya que el movimiento del tronco invalida la posición lateral verdadera. La descripción ``como si quisieras mirar el techo detrás tuyo'' resulta útil para que el paciente comprenda el grado de extensión requerido.
	
	
	\onecolumn
	\clearpage
	
	\begin{figure}
		\centering
		\includegraphics[width=0.7\linewidth]{imagenes/col13}
		\caption{}
		\label{fig:col13}
	\end{figure}
	
	\twocolumn
	\clearpage
	%----------------------------------------------------------
	%  Para agregar otro enfoque duplicá el bloque
	%  desde \subsection hasta \textbf{Trucos}
	%----------------------------------------------------------
	
	
%%%%%%%%%%%%%%%%%%%%%%%%%%%%%%%%%%%%%%%%%%%%%%%%%%%%%%%%%%%
%  SECCIÓN: Columna Cervical
%  Archivo: secciones/columna_cervical.tex
%  Sin preámbulo ni \begin{document} — lo maneja main.tex
	%%%%%%%%%%%%%%%%%%%%%%%%%%%%%%%%%%%%%%%%%%%%%%%%%%%%%%%%%%%
	
	\section{Columna Cervical}
	
	\subsection{Oblicua Posterior Derecha (OPD) - Oblicua Posterior Izquierda (ODI) - AP}
	
	\subsubsection*{1. Identificación}
	
	\textbf{Tipo:} Proyección complementaria. Forma par con la proyección oblicua contralateral (OPD con ODI y viceversa); siempre se realizan en conjunto.
	
	\textbf{Indicaciones clínicas principales:}
	\begin{itemize}
		\item Estenosis cervical de los orificios intervertebrales (compresión del canal espinal y raíces nerviosas)
		\item Osteofitosis y discopatías cervicales
		\item Cervicobraquialgia, radiculopatía y parestesias (dolor al flexionar el cuello o adormecimiento de manos)
	\end{itemize}
	
	\subsubsection*{2. Técnica Base}
	
	\textbf{Receptor de imagen:} 18$\times$24 cm o 24$\times$30 cm. Orientación \textbf{longitudinal} (vertical). El RI se centra a nivel de C3, aproximadamente 2{,}5 cm por encima del punto más prominente del cartílago tiroides, para compensar la angulación craneal del rayo central.
	
	\textbf{DFRI:} 150--180 cm. La distancia mayor se emplea para compensar el aumento de OID generado por la presencia de los hombros respecto al receptor de imagen.
	
	\textbf{Técnica:} Homogeneizada (escala larga, kVp suficiente). Permite atravesar y visualizar los diferentes espesores desiguales presentes en la proyección: hombro, hueso, aire y tejido blando cervical.
	
	\textbf{Bucky:} Bucky mural o de mesa con grilla, según la posición del paciente.
	
	\subsubsection*{3. Posicionamiento}
	
	\paragraph{Paciente}
	De pie (preferido por comodidad), sentado o en decúbito semisupino con la espalda apoyada contra el receptor. Se parte de la posición AP estándar con el plano medio sagital perpendicular al RI.
	
	\begin{figure}[h]
		\centering
		\includegraphics[width=0.7\linewidth]{imagenes/col14}
		\caption{Proyección AP axial oblicua en bipedestación de los agujeros intervertebrales derechos: posición\textbf{ OPI}.}
		\label{fig:col14}
	\end{figure}
	
	
	\paragraph{Región de interés}
	El cuerpo y la cabeza se rotan conjuntamente 45° respecto al plano del RI, apoyando la espalda contra el receptor. Se solicita al paciente que eleve y adelante el mentón para evitar superposición con las vértebras superiores, y que baje los hombros para no cubrir C7. Es fundamental que el mentón no gire hacia un lado, para evitar rotación de las primeras vértebras cervicales superiores. En posición semisupino: se colocan apoyos bajo la parte inferior del tórax y la cadera elevada, y un apoyo adicional bajo la cabeza.
	
	\paragraph{Rayo Central}
	Angulado 15--20° en dirección \textbf{craneal}, dirigido a \textbf{C4 (nuez de Adán)}. La angulación craneal responde a la orientación caudocraneal de los agujeros intervertebrales (inclinados aproximadamente 15° de atrás hacia adelante), permitiendo un paso tangencial a través de ellos y evitando que los hombros cubran C7.
	
	\paragraph{Receptor de imagen}
	Centrado a nivel de C3 (2{,}5 cm por encima del cartílago tiroides más prominente). Colimación a los bordes laterales de las partes blandas del cuello y a los bordes superior e inferior del RI. La región a incluir abarca desde la base del cráneo hasta los cuerpos vertebrales C1--C7 y T1, visualizando la transición cervicotorácica.
	
	\paragraph{Respiración}
	Suspendida al momento de la exposición. Permite eliminar el movimiento de las estructuras cervicales y garantiza nitidez en la imagen.
	
	\subsubsection*{4. Anatomía Visible}
	
	Se demuestran los agujeros intervertebrales del lado \textbf{más alejado del receptor} (C2--C3 a C7--T1) en posición abierta, permitiendo evaluar estenosis foraminal. Los pedículos del lado más alejado del RI quedan proyectados en visión de perfil. Los espacios discales intervertebrales aparecen abiertos por la curvatura cervical. Los cuerpos vertebrales se proyectan en oblicuo. El hueso occipital puede superponerse levemente a C1, lo cual es esperado y aceptable. Los hombros no deben cubrir C7.
	
	\subsubsection*{5. Criterios de Evaluación}
	
	\textbf{Posicionamiento correcto:} Agujeros intervertebrales del lado más alejado del receptor abiertos y de tamaño y contorno uniformes desde C2--C3 hasta C7--T1. Espacios discales intervertebrales abiertos. Barbilla elevada, sin superposición de la mandíbula sobre las vértebras superiores (atlas y axis). Hueso occipital superpuesto a C1 (aceptable y esperado). Hombros sin cubrir C7.
	
	\textbf{Ausencia de rotación:} Rotación homogénea de 45° del tronco y la cabeza. Si los agujeros intervertebrales aparecen cerrados, indica corrección insuficiente de la rotación del tronco. Tamaño y contorno uniformes de todos los agujeros visibles.
	
	\textbf{Penetración y contraste:} Técnica homogeneizada con escala larga (alto kVp). Debe permitir diferenciar simultáneamente los distintos espesores: tejido blando cervical, hueso, aire y región de hombros.
	
	\subsubsection*{6. Puntos Críticos}
	
	\textbf{Agujeros intervertebrales cerrados o vértebras amontonadas:} Indica insuficiente angulación craneal del RC o rotación de tronco incorrecta. Corregir aumentando la angulación craneal (hasta 20°) y verificando los 45° de rotación del cuerpo.
	
	\textbf{Mandíbula superpuesta a atlas y axis:} El mentón no fue elevado o adelantado correctamente. Solicitar al paciente que eleve y protruya el mentón. Verificar que no gire el mentón lateralmente para evitar rotación de las vértebras superiores.
	
	\textbf{Adaptaciones:} En paciente con trauma o limitación de movimiento, puede realizarse en decúbito semisupino con apoyos bajo tórax y cadera. La posición supina es equivalente en resultado si la rotación de 45° se mantiene correcta.
	
	\textbf{Trucos:} Las oblicuas cervicales posteriores se marcan como proyecciones de frente (no de perfil): la marca derecha corresponde al lado derecho de la imagen. Para identificar el lado en la radiografía, observar la posición de la mandíbula: si queda hacia la izquierda, corresponde a oblicua posterior derecha o anterior izquierda. Siempre se realizan ambas oblicuas (OPD y ODI) para comparación bilateral.
	
	
	\onecolumn
	\clearpage
	
	\begin{figure}
		\centering
		\includegraphics[width=0.7\linewidth]{imagenes/col16}
		\caption{Posición OPD para la demostración del lado izquierdo.
		}
		\label{fig:col16}
	\end{figure}
		
	
	\twocolumn
	\clearpage
	
	\onecolumn
	\clearpage
	
	\begin{figure}
		\centering
		\includegraphics[width=0.7\linewidth]{imagenes/col15}
		\caption{Posición OPI para la demostración del lado derecho.
		}
		\label{fig:col15}
	\end{figure}
		
	
	\twocolumn
	\clearpage
	%----------------------------------------------------------
	%  Para agregar otro enfoque duplicá el bloque
	%  desde \subsection hasta \textbf{Trucos}
	%----------------------------------------------------------
	
	
	%%%%%%%%%%%%%%%%%%%%%%%%%%%%%%%%%%%%%%%%%%%%%%%%%%%%%%%%%%%
	%  SECCIÓN: Columna Cervical — PA Oblicuas Anteriores
	%  Archivo: secciones/oblicuas_anteriores_cervical.tex
	%  Sin preámbulo ni \begin{document} — lo maneja main.tex
		%%%%%%%%%%%%%%%%%%%%%%%%%%%%%%%%%%%%%%%%%%%%%%%%%%%%%%%%%%%
		
		\section{Columna Cervical}
		
		\subsection{Oblicua Anterior Derecha (OAD) - Oblicua Anterior Izquierda (OAI)}
		
		\subsubsection*{1. Identificación}
		
		\textbf{Tipo:} Complementaria. Forma par con las oblicuas posteriores (OPD/OPI); ambos lados se realizan siempre para comparación bilateral.
		
		\textbf{Indicaciones clínicas principales:}
		\begin{itemize}
			\item Estenosis de los orificios intervertebrales cervicales (compresión del canal espinal y raíces nerviosas)
			\item Osteofitosis y discopatías cervicales
			\item Cervicobraquialgia, radiculopatía y parestesias (dolor al flexionar el cuello o adormecimiento de manos)
		\end{itemize}
		
		\subsubsection*{2. Técnica Base}
		
		\textbf{Receptor de imagen:} 18$\times$24 cm o 24$\times$30 cm. Orientación \textbf{longitudinal}. Centrado a la altura de C5 (2,5 cm caudal al punto más prominente del cartílago tiroides).
		
		\textbf{DFRI:} 150--180 cm. La mayor distancia compensa el aumento de OID propio de la proyección oblicua y mejora la nitidez de la imagen.
		
		\textbf{Técnica:} Homogeneizada (escala larga, kVp suficiente). Permite atravesar y visualizar los diferentes espesores de la región: hombro, hueso, aire y tejido blando cervical.
		
		\textbf{Bucky:} Bucky de pared con grilla.
		
		\subsubsection*{3. Posicionamiento}
		
		\paragraph{Paciente}
		De pie o sentado en bipedestación, con el hombro apoyado contra el Bucky de pared. El paciente queda de cara al tubo (posición PA). También puede realizarse en decúbito prono (semiprono). Los brazos se mantienen a los lados del cuerpo; en espiración los hombros descienden, mejorando la visualización de C7.
		
		\begin{figure}[h]
			\centering
			\includegraphics[width=0.7\linewidth]{imagenes/col17}
			\caption{}
			\label{fig:col17}
		\end{figure}
		
		\begin{figure}[h]
			\centering
			\includegraphics[width=0.7\linewidth]{imagenes/col18}
			\caption{}
			\label{fig:col18}
		\end{figure}
		
		
		
		\paragraph{Región de interés}
		Se rota todo el cuerpo del paciente 45° respecto al receptor de imagen, de modo que los agujeros intervertebrales del lado más alejado del RI queden en paralelo al plano del receptor. El plano medio sagital debe quedar alineado con el eje de la columna. Se eleva y hace protruir la barbilla lo suficiente para evitar la superposición de la mandíbula sobre las vértebras cervicales superiores (atlas y axis). En posición semiprono se utiliza el antebrazo y la rodilla flexionada del lado elevado como apoyo.
		
		\paragraph{Rayo Central}
		Angulado \textbf{15--20° en dirección caudal}, dirigido hacia \textbf{C4 (nuez de Adán)}. La angulación caudal coincide con la dirección descendente de los agujeros intervertebrales vistos desde atrás, permitiendo proyectarlos abiertos.
		
		\paragraph{Receptor de imagen}
		Centrado a la altura de C5. La colimación abarca los bordes laterales de las partes blandas del cuello y los bordes superior e inferior coinciden con los bordes del RI. La región a incluir es desde la base del cráneo hasta T1.
		
		\paragraph{Respiración}
		Suspendida en espiración. La espiración favorece el descenso de los hombros, liberando la visualización de C7 y reduciendo la necesidad de aumentar la angulación.
		
		\subsubsection*{4. Anatomía Visible}
		
		Agujeros intervertebrales del lado \textbf{más alejado del RI} (los más próximos al tubo), desde C2--C3 hasta C7--T1, proyectados abiertos. Pedículos del lado más cercano al receptor. Cuerpos vertebrales cervicales en proyección oblicua con espacios interdiscales abiertos. Articulaciones uncovertebrales. La mandíbula no debe superponerse al atlas ni al axis; el hueso occipital no debe superponerse al axis; los hombros no deben tapar C7.
		
		\subsubsection*{5. Criterios de Evaluación}
		
		\textbf{Posicionamiento correcto:} Agujeros intervertebrales del lado alejado al RI abiertos y de tamaño y contorno uniformes desde C2 hasta C7--T1. Espacios interdiscales abiertos. Barbilla elevada sin superposición mandibular sobre vértebras superiores. Todas las vértebras cervicales incluidas hasta T1.
		
		\textbf{Ausencia de rotación:} Rotación corporal exacta de 45°; si los agujeros aparecen cerrados, la rotación del tronco es insuficiente o excesiva y debe corregirse.
		
		\textbf{Penetración y contraste:} Técnica homogeneizada que permite diferenciar simultáneamente el hombro, el hueso, el aire y el tejido blando cervical en una escala de grises larga.
		
		\subsubsection*{6. Puntos Críticos}
		
		\textbf{Agujeros intervertebrales cerrados:} Indica rotación corporal inadecuada (diferente a 45°). Corregir la rotación del tronco hasta lograr el ángulo exacto.
		
		\textbf{Superposición mandibular sobre vértebras superiores:} La barbilla no fue elevada ni protraída suficientemente. Indicar al paciente que eleve y adelante el mentón.
		
		\textbf{Adaptaciones:} Si los hombros tapan C7, aumentar la angulación caudal o indicar al paciente que baje activamente los hombros. En pacientes con limitación de movilidad puede realizarse en decúbito semiprono usando antebrazo y rodilla como apoyo.
		
		\textbf{Trucos:} La angulación es caudal (opuesta a las oblicuas posteriores AP). Se prefieren las oblicuas anteriores sobre las posteriores por la menor dosis de radiación a la tiroides (tiroides queda en RP). Para identificar la proyección en la radiografía: si el mentón apunta hacia la marca derecha es OAI u OPD; si apunta hacia la marca izquierda es OAD u OPI. Las oblicuas anteriores demuestran los agujeros del lado \textbf{más alejado} del receptor (a diferencia de las posteriores, que muestran los más cercanos).
		
		\onecolumn
		\clearpage
		
		\begin{figure}
			\centering
			\includegraphics[width=0.7\linewidth]{imagenes/col20}
			\caption{Posición OAI para la demostración del lado izquierdo.
			}
			\label{fig:col20}
		\end{figure}
		
		
		\twocolumn
		\clearpage
		
		
		\onecolumn
		\clearpage
		
		\begin{figure}
			\centering
			\includegraphics[width=0.7\linewidth]{imagenes/col19}
			\caption{Posición OAD para la demostración del lado derecho.}
			\label{fig:col19}
		\end{figure}
		
		
		\twocolumn
		\clearpage
		
		
		%----------------------------------------------------------
		%  Para agregar otro enfoque duplicá el bloque
		%  desde \subsection hasta \textbf{Trucos}
		%----------------------------------------------------------
		
		
		
	%%%%%%%%%%%%%%%%%%%%%%%%%%%%%%%%%%%%%%%%%%%%%%%%%%%%%%%%%%%
	%  SECCIÓN: Columna Torácica
	%  Archivo: secciones/pasaje_cervico_dorsal.tex
	%  Sin preámbulo ni \begin{document} — lo maneja main.tex
		%%%%%%%%%%%%%%%%%%%%%%%%%%%%%%%%%%%%%%%%%%%%%%%%%%%%%%%%%%%
		
		\section{Columna Torácica}
		
		\subsection{Pasaje Cérvico-Dorsal - Técnica del Nadador (Método de Twining)}
		
		\subsubsection*{1. Identificación}
		
		\textbf{Tipo:} Proyección complementaria. Forma par con la proyección lateral de columna torácica.
		
		\textbf{Indicaciones clínicas principales:}
		\begin{itemize}
			\item Traumatismo cervical, fractura o luxación de C7 o T1.
			\item Cervicobraquialgia: dolor irradiado a los brazos y la base del cuello.
			\item Cuello rígido.
			\item Proyección de elección cuando los hombros prominentes ocultan C7--T1 en la lateral cervical estándar, impidiendo la visualización de la cabeza del húmero.
			\item Proyección lateral de vértebras torácicas superiores y región de transición C7--T1.
		\end{itemize}
		
		\subsubsection*{2. Técnica Base}
		
		\textbf{Receptor de imagen:} 24x30 cm. Orientación \textbf{longitudinal} (vertical), centrado a la altura del espacio C7--T1, localizado 5 cm por encima de la escotadura yugular. \textbf{Filtro de compensación obligatorio:} las diferencias extremas de densidad entre el cuello fino y la región torácica gruesa exigen el uso de un filtro especialmente diseñado para demostrar el área C7--T1 en una sola imagen.
		
		\textbf{DFRI:} 150 o 180 cm. La mayor distancia reduce la penumbra geométrica y minimiza la magnificación de estructuras profundas en esta región anatómicamente compleja.
		
		\textbf{Técnica:} Intermedia. Se utiliza un kVp suficiente para atravesar el hombro y un mAs elevado para ganar densidad en los cuerpos vertebrales sin sobreexponer el aire del cuello. Los valores dependen del biotipo del paciente. Técnica respiratoria opcional: bajo mAs con tiempo de exposición de 3--4 segundos, solicitando al paciente respiraciones breves y regulares para difuminar la anatomía pulmonar y mejorar la visualización ósea.
		
		\textbf{Bucky:} Sí, con grilla. El uso de grilla es necesario dada la cantidad de tejido atravesado en la región del hombro y la parte superior del tórax.
		
		\subsubsection*{3. Posicionamiento}
		
		\paragraph{Paciente}
		El paciente se coloca en posición lateral verdadera, de pie (preferible), sentado o en decúbito lateral. El plano medio sagital debe quedar paralelo al plano del RI. En decúbito lateral, la cabeza se eleva sobre el brazo del paciente o un soporte firme; puede colocarse un apoyo bajo la parte inferior del tórax para alinear la columna. Si el paciente está incorporado, se apoya en un dispositivo de rejilla vertical.
		
		\begin{figure}[h]
			\centering
			\includegraphics[width=0.7\linewidth]{imagenes/col21}
			\caption{}
			\label{fig:col21}
		\end{figure}
		
		\begin{figure}[h]
			\centering
			\includegraphics[width=0.7\linewidth]{imagenes/col22}
			\caption{}
			\label{fig:col22}
		\end{figure}
		
		
		
		\paragraph{Región de interés}
		El plano medio coronal del cuerpo se centra en la línea media de la rejilla. El brazo próximo al RI se extiende por encima de la cabeza desplazando la cabeza humeral hacia anterior (recomendado) para que no se superponga a la columna dorsal. Si el paciente está incorporado, el codo se flexiona y el antebrazo se apoya sobre la cabeza. El brazo alejado del RI se desciende lo máximo posible a lo largo del cuerpo, traccionándolo hacia posterior para separar ambos hombros en sentido vertical. Con estos dos movimientos se liberan las cabezas humerales y la columna cérvico-dorsal queda sin superposición. Se recomienda una leve rotación del paciente hacia anterior junto con el brazo elevado para terminar de separar las estructuras.
		
		\paragraph{Rayo Central}
		Perpendicular al plano del receptor, dirigido al espacio C7--T1, localizado \textbf{5 cm por encima de la escotadura yugular}. Si el paciente no puede descender el hombro alejado del RI, se aplica una angulación \textbf{caudal de 3--5°} para compensar la diferencia de altura entre los hombros.
		
		\paragraph{Receptor de imagen}
		Centrado a la altura del espacio C7--T1. La región de inclusión debe abarcar superiormente desde C4--C5 e inferiormente hasta T3--T4 para asegurar la visualización completa de la transición cérvico-torácica, incluyendo cuerpos vertebrales y apófisis espinosas. Siempre debe visualizarse el pasaje de una región a otra (C7--T1).
		
		\paragraph{Respiración}
		Suspendida en espiración. Si el paciente coopera, puede utilizarse la \textbf{técnica respiratoria opcional}: se solicitan respiraciones breves y regulares durante un tiempo de exposición de 3--4 segundos con bajo mAs, lo que difumina la anatomía pulmonar y resalta las estructuras óseas vertebrales.
		
		\subsubsection*{4. Anatomía Visible}
		
		Se visualizan las vértebras cervicales inferiores y las vértebras torácicas superiores entre los hombros. La región de transición cérvico-torácica (C7--T1) debe quedar claramente demostrada. Los cuerpos vertebrales y las apófisis espinosas deben incluirse en el campo. La separación de las cabezas humerales mediante la posición diferencial de los brazos permite liberar la columna de superposiciones. Los agujeros intervertebrales torácicos se visualizan en perfil sin necesidad de enfoques específicos, dado que se ubican a 45° respecto al plano sagital y en posición lateral ya se proyectan adecuadamente.
		
		\subsubsection*{5. Criterios de Evaluación}
		
		\textbf{Posicionamiento correcto:} Vértebras cervicales inferiores y torácicas superiores visibles entre los hombros, sin rotación significativa respecto a una posición lateral verdadera. El espacio C7--T1 debe estar centrado en la imagen. La región incluye desde C4--C5 hasta T3--T4, con cuerpos vertebrales y apófisis espinosas presentes.
		
		\textbf{Ausencia de rotación:} Cuerpos vertebrales de perfil, sin doble contorno. Si presentan doble contorno, el PMS no está paralelo al RI y debe corregirse. Las cabezas humerales deben estar mínimamente superpuestas a la columna vertebral. En posición lateral estricta la cabeza humeral tiende a superponerse al pasaje cérvico-dorsal; debe evitarse mediante la separación diferencial de los brazos.
		
		\textbf{Penetración y contraste:} Adecuada penetración de los rayos X a través de la región del hombro, con densidad suficiente para visualizar los cuerpos vertebrales sin sobreexponer el aire del cuello. El filtro de compensación debe haber equilibrado la diferencia de densidades entre el cuello y la región torácica superior, produciendo una imagen de densidad uniforme en toda la columna estudiada.
		
		\subsubsection*{6. Puntos Críticos}
		
		\textbf{Superposición humeral sobre la columna:} Si las cabezas humerales se superponen al espacio C7--T1, la proyección pierde su utilidad diagnóstica. Corrección: asegurar la máxima separación vertical de los hombros mediante la extensión del brazo proximal sobre la cabeza y el descenso del brazo distal; aplicar leve rotación anterior del brazo elevado.
		
		\textbf{Diferencia extrema de densidades:} Sin filtro de compensación, el cuello queda sobreexpuesto y la región torácica subexpuesta en la misma imagen, impidiendo la evaluación de C7--T1. Corrección: el uso del filtro de compensación es obligatorio en esta proyección.
		
		\textbf{Adaptaciones:} Si el paciente no puede descender el hombro alejado del RI, angular el RC caudal 3--5° incidiendo en C7--T1. Si en la imagen obtenida la cabeza humeral queda superpuesta al pasaje cérvico-dorsal por no haber oblicuado el brazo elevado ni movilizado el hombro contralateral, repetir la proyección corrigiendo el posicionamiento o angulando el RC. En decúbito lateral, colocar apoyo bajo el tórax inferior para alinear la columna con el RI.
		
		\textbf{Trucos:} La técnica respiratoria (3--4 segundos, respiraciones breves regulares) es especialmente útil para difuminar estructuras pulmonares y mejorar la nitidez ósea vertebral. La recomendación de desplazar el húmero del brazo elevado hacia anterior evita que se proyecte sobre la columna dorsal. Siempre verificar que el pasaje C7--T1 quede incluido y visible, ya que es el objetivo diagnóstico central de esta proyección.
		
		
		
		\clearpage
		\onecolumn
		
		\begin{figure}
			\centering
			\includegraphics[width=0.7\linewidth]{imagenes/col23}
			\caption{}
			\label{fig:col23}
		\end{figure}
		
		
		\twocolumn
		\clearpage
		%----------------------------------------------------------
		%  Para agregar otro enfoque duplicá el bloque
		%  desde \subsection hasta \textbf{Trucos}
		%----------------------------------------------------------
		
		
	%%%%%%%%%%%%%%%%%%%%%%%%%%%%%%%%%%%%%%%%%%%%%%%%%%%%%%%%%%%
	%  SECCIÓN: Columna Dorsal
	%  Archivo: secciones/columna_dorsal.tex
	%  Sin preámbulo ni \begin{document} — lo maneja main.tex
		%%%%%%%%%%%%%%%%%%%%%%%%%%%%%%%%%%%%%%%%%%%%%%%%%%%%%%%%%%%
		
		\section{Columna Torácica}
		
		\subsection{Frente de Columna Torácica}
		
		\subsubsection*{1. Identificación}
		
		\textbf{Tipo:} Proyección básica. Forma par con la proyección lateral de columna dorsal.
		
		\textbf{Indicaciones clínicas principales:}
		\begin{itemize}
			\item Dorsalgia crónica
			\item Cifosis (joroba)
			\item Fracturas por compresión
			\item Escoliosis
		\end{itemize}
		
		\subsubsection*{2. Técnica Base}
		
		\textbf{Receptor de imagen:} 35$\times$43 cm. Orientación \textbf{longitudinal}. Borde superior del RI ubicado 4 a 5 cm por encima de los hombros, de forma que se incluya el pasaje C7-T1 en el extremo superior y el pasaje T12-L1 en el extremo inferior, permitiendo contar e identificar cada vértebra.
		
		\textbf{DFRI:} 100 cm.
		
		\textbf{Técnica:} Homogeneizada (alto kVp). Se utiliza alto kVp para engreír las distintas densidades y poder atravesar todo el espesor de la región, compensando la diferencia de grosor entre el extremo superior (tórax) y el extremo inferior (abdomen superior). Se orienta el ánodo hacia la cabeza y el cátodo hacia los pies, aprovechando que el cátodo emite mayor intensidad de rayos para atravesar la zona de mayor densidad.
		
		\textbf{Bucky:} Bucky de mesa o de pared. Grilla sí.
		
		\subsubsection*{3. Posicionamiento}
		
		\paragraph{Paciente}
		Decúbito supino (sin almohada para no acentuar la cifosis torácica) o en bipedestación. El plano medio sagital debe quedar perpendicular al receptor de imagen. En posición supina, se flexionan caderas y rodillas con los muslos en posición vertical para reducir la cifosis torácica y lograr un contacto más uniforme de la columna con el receptor. Los brazos se disponen a los lados del cuerpo y los hombros deben encontrarse en un mismo plano.
		
		\begin{figure}[h]
			\centering
			\includegraphics[width=0.7\linewidth]{imagenes/col25}
			\caption{}
			\label{fig:col25}
		\end{figure}
		
		
		\paragraph{Región de interés}
		Columna alineada y centrada en la porción media del receptor. Los hombros deben estar en un mismo plano, sin rotación del tronco, de modo que el plano medio sagital quede perpendicular al RI.
		
		\paragraph{Rayo Central}
		Perpendicular al plano del receptor. Punto de entrada: \textbf{T6}, que se localiza en el punto medio entre la escotadura yugular (T3) y el apéndice xifoides (T9), aproximadamente una palma de ancho hacia cada lado. Colimar estrechamente sobre la columna torácica.
		
		\paragraph{Receptor de imagen}
		Centrado en D6. Borde superior 4 a 5 cm por encima de los hombros para asegurar la inclusión del pasaje C7-D1 superiormente y del pasaje T12-L1 inferiormente.
		
		\paragraph{Respiración}
		Suspendida al final de una espiración forzada. Al espirar, el diafragma asciende, logrando una densidad más uniforme que facilita la visualización de las vértebras inferiores (T9-T12), que de otro modo quedarían ocultas por la sombra diafragmática.
		
		\subsubsection*{4. Anatomía Visible}
		
		Se deben visualizar las 12 vértebras torácicas en su totalidad, incluyendo los cuerpos vertebrales, apófisis espinosas, pedículos y costillas articuladas. Se debe incluir el pasaje cérvico-dorsal (C7-T1) y el pasaje dorso-lumbar (T12-L1), imprescindibles para la identificación y conteo correcto de cada vértebra. Con receptor de 35$\times$43 cm también deben visualizarse las columnas laterales, los hombros, los campos pulmonares y el diafragma. La identificación se realiza buscando la primera vértebra que porta costilla (T1) desde arriba y verificando que la última vértebra visible también las posea (T12); puede haber excepciones por costillas o vértebras supernumerarias, que se identifican por ser de menor tamaño o no coincidir con el conteo habitual.
		
		\subsubsection*{5. Criterios de Evaluación}
		
		\textbf{Posicionamiento correcto:} Las 12 vértebras torácicas deben estar visibles con densidad uniforme. Si no es posible lograr densidad uniforme en una sola toma, se obtienen dos radiografías separadas: una para vértebras superiores y otra para las inferiores. La columna vertebral debe estar alineada en la línea media de la radiografía. El haz debe estar colimado sobre la columna torácica.
		
		\textbf{Ausencia de rotación:} Las articulaciones esterno claviculares deben estar equidistantes respecto a la línea media de la columna. Las apófisis espinosas deben estar centradas en los cuerpos vertebrales. Si las espinosas aparecen desviadas, el paciente está rotado: corregir colocando el plano medio sagital perpendicular al RI.
		
		\textbf{Penetración y contraste:} La técnica homogeneizada de alto kVp debe mostrar densidad uniforme a lo largo de toda la columna, permitiendo visualizar tanto las vértebras superiores (menor espesor) como las inferiores (mayor espesor abdominal) en una sola imagen.
		
		\subsubsection*{6. Puntos Críticos}
		
		\textbf{Rotación del paciente:} Las apófisis espinosas aparecen desviadas del centro de los cuerpos vertebrales y las articulaciones esterno claviculares pierden su equidistancia. Corregir asegurando que el plano medio sagital quede estrictamente perpendicular al RI.
		
		\textbf{Cifosis no compensada:} Si no se flexionan caderas y rodillas en decúbito supino, la cifosis torácica acentuada aleja parte de la columna del receptor, aumentando el OID y degradando la nitidez. Corregir flexionando caderas y rodillas o utilizando la posición en bipedestación.
		
		\textbf{Adaptaciones:} En pacientes con cifosis severa o en trauma, puede utilizarse la posición en bipedestación o con apoyo adicional. En pacientes con diferencia de grosor muy marcada puede requerirse el uso de filtro de compensación o dos exposiciones separadas.
		
		\textbf{Trucos:} Orientar el cátodo hacia los pies (zona de mayor densidad del abdomen superior) aprovecha el efecto de talón del ánodo para lograr mayor intensidad donde más se necesita. La espiración forzada eleva el diafragma y homogeneiza la densidad de las vértebras inferiores. Para contar vértebras, identificar T1 como la primera con costilla desde arriba y verificar que T12 sea la última con costilla hacia abajo.
		
		
		\onecolumn
		\clearpage
		
		\begin{figure}
			\centering
			\includegraphics[width=0.7\linewidth]{imagenes/col24}
			\caption{El extremo catódico (-) del tubo de rayos X sobre la parte inferior del tórax (una densidad más uniforme).
			}
			\label{fig:col24}
		\end{figure}
		
		
		\twocolumn
		\clearpage
		
		%----------------------------------------------------------
		%  Para agregar otro enfoque duplicá el bloque
		%  desde \subsection hasta \textbf{Trucos}
		%----------------------------------------------------------
		
		
	%%%%%%%%%%%%%%%%%%%%%%%%%%%%%%%%%%%%%%%%%%%%%%%%%%%%%%%%%%%
	%  SECCIÓN: Columna Torácica
	%  Archivo: secciones/columna_toracica.tex
	%  Sin preámbulo ni \begin{document} — lo maneja main.tex
		%%%%%%%%%%%%%%%%%%%%%%%%%%%%%%%%%%%%%%%%%%%%%%%%%%%%%%%%%%%
		
		\section{Columna Torácica}
		
		\subsection{Lateral de Columna Torácica (Perfil)}
		
		\subsubsection*{1. Identificación}
		
		\textbf{Tipo:} Proyección complementaria. Forma par con la proyección AP de columna torácica.
		
		\textbf{Indicaciones clínicas principales:}
		\begin{itemize}
			\item Fracturas por compresión vertebral
			\item Subluxación vertebral
			\item Cifosis
		\end{itemize}
		
		\subsubsection*{2. Técnica Base}
		
		\textbf{Receptor de imagen:} 35$\times$43 cm. Orientación \textbf{longitudinal}. El borde superior del RI se ubica 4--5 cm por encima de los hombros.
		
		\textbf{DFRI:} 100 cm.
		
		\textbf{Técnica:} Técnica homogenizada (técnica respiratoria). Se separa el miliamperaje de los segundos; se aumenta el tiempo de exposición a 1400--1600 ms (equivalente a un ciclo respiratorio completo). Durante el disparo el paciente respira con normalidad (inspira y espira), lo que genera borrosidad cinética de los arcos costales, dejando la columna con mayor nitidez. La técnica respiratoria implica mayor dosis, por lo que no se indica en niños ni en pacientes con dolor agudo intenso.
		
		\textbf{Bucky:} Bucky de pared o de mesa según posición del paciente. Grilla obligatoria dada la técnica utilizada.
		
		\subsubsection*{3. Posicionamiento}
		
		\paragraph{Paciente}
		Decúbito lateral izquierdo (para que el corazón quede más próximo al RI) o en bipedestación. El paciente viste bata abierta por detrás de la columna. Se coloca almohada bajo la cabeza para mantener el eje longitudinal de la columna vertebral en posición horizontal. Se flexionan caderas y rodillas. El plano medio sagital debe mantenerse paralelo al RI.
		
		\begin{figure}[h]
			\centering
			\includegraphics[width=0.7\linewidth]{imagenes/col27}
			\caption{}
			\label{fig:col27}
		\end{figure}
		
		\begin{figure}[h]
			\centering
			\includegraphics[width=0.7\linewidth]{imagenes/col26}
			\caption{}
			\label{fig:col26}
		\end{figure}
		
		
		\paragraph{Región de interés}
		La región abarca desde C7 hasta L1--L2, incluyendo los pasajes cervicotorácico y toracolumbar. Los brazos del paciente se dirigen hacia adelante (no hacia arriba), con el objetivo de desplazar el húmero fuera del pasaje cervicotorácico, desplegar los agujeros intervertebrales y elevar las costillas. Elevar los brazos provocaría superposición del húmero sobre T1--T3. Si la columna no es paralela al RI, se colocan soportes bajo la zona lumbar y se alinean hombros y cadera. Cuando el paciente está en decúbito, la espalda debe quedar orientada hacia el operador.
		
		\paragraph{Rayo Central}
		Perpendicular al plano del receptor. Punto de entrada: \textbf{T6/T7 (borde inferior de la escápula)}, cuatro traveses de dedos por dentro del reborde cutáneo posterior.
		
		\paragraph{Receptor de imagen}
		Centrado en T6/T7. Borde superior del RI ubicado 4--5 cm por encima de los hombros, de manera que el campo abarque desde C7 hasta L1--L2.
		
		\paragraph{Respiración}
		Técnica especial de respiración: se le indica al paciente que respire normalmente (inspire y espire) durante todo el tiempo de exposición. El tiempo de disparo se extiende a 1400--1600 ms para abarcar al menos un ciclo respiratorio completo. El movimiento de los arcos costales genera borrosidad cinética sobre éstos, eliminando su superposición sobre los cuerpos vertebrales y permitiendo visualizar la columna con mayor nitidez. Esto se obtiene separando el mA de los segundos en el equipo y ajustando el tiempo de exposición según el foco utilizado.
		
		\subsubsection*{4. Anatomía Visible}
		
		Se visualizan las 12 vértebras torácicas de perfil, con cuerpos vertebrales superpuestos entre sí (imagen de cuerpo único) y espacios discales abiertos. La región abarca desde C7 hasta L1--L2. Los arcos costales aparecen borrados por borrosidad cinética (técnica respiratoria). Las vértebras T1 a T3 pueden no visualizarse correctamente por superposición de los hombros. La línea posterolateral de los bordes posteriores debe presentarse como una línea única continua, lo que confirma la correcta superposición de los cuerpos vertebrales.
		
		\subsubsection*{5. Criterios de Evaluación}
		
		\textbf{Posicionamiento correcto:} Cuerpos vertebrales superpuestos de perfil formando una imagen única; espacios discales abiertos en toda la extensión; las 12 vértebras torácicas centradas en el RI; bordes posteriores formando una línea posterolateral única y continua.
		
		\textbf{Ausencia de rotación:} Superposición simétrica de las costillas homólogas entre sí; ausencia de doble línea en los bordes posteriores de los cuerpos vertebrales. La presencia de doble línea o falta de superposición costal indica rotación del paciente.
		
		\textbf{Penetración y contraste:} Amplia latitud de exposición. Vértebras claramente visibles a través de las sombras costales y pulmonares. Los arcos costales presentan borrosidad cinética uniforme. Colimación estrecha para reducir la radiación dispersa.
		
		\subsubsection*{6. Puntos Críticos}
		
		\textbf{Rotación del paciente:} Genera doble línea en los bordes posteriores de los cuerpos vertebrales y pérdida de superposición costal. Corrección: alinear hombros y caderas en el mismo plano, colocar soportes bajo la zona lumbar para que el eje longitudinal de la columna quede paralelo al RI, flexionar rodillas para estabilizar la posición.
		
		\textbf{No visualización de T1--T3:} Producida por superposición de los hombros y subexposición relativa en esa zona. El número de vértebras visibles depende del tamaño y la forma del paciente. No es un error corregible en todos los casos; se documenta como limitación de la proyección.
		
		\textbf{Adaptaciones:} En pacientes con dolor agudo o pediátricos se evalúa no utilizar la técnica respiratoria, dado que implica mayor tiempo de exposición y mayor dosis. En decúbito lateral, si la columna no es paralela al RI, se colocan soportes bajo la región lumbar para corregir la angulación.
		
		\textbf{Trucos:} La posición en decúbito lateral izquierdo se prefiere para acercar el corazón al RI y reducir la magnificación cardíaca. Los brazos deben dirigirse hacia adelante y no hacia arriba: esta diferencia es clave para evitar superposición del húmero en el pasaje cervicotorácico y para elevar las costillas. Al separar mA de segundos en el equipo, cambiar el foco recalcula automáticamente los milisegundos necesarios. La técnica respiratoria equivale a mayor dosis: la indicación depende de la calidad diagnóstica requerida y del estado del paciente.
		
		
		\onecolumn
		\clearpage
		
		\begin{figure}
			\centering
			\includegraphics[width=0.7\linewidth]{imagenes/col28}
			\caption{Proyección lateral torácica.\textbf{ A}. Respiración suspendida con una exposición de 0,5 s. \textbf{B}. Técnica en respiración con una exposición de 1,5 s. Obsérvese que las marcas pulmonares están difuminadas.
			}
			\label{fig:col28}
		\end{figure}
		
		
		\twocolumn
		\clearpage
		
		%----------------------------------------------------------
		%  Para agregar otro enfoque duplicá el bloque
		%  desde \subsection hasta \textbf{Trucos}
		%----------------------------------------------------------
		
		
	%%%%%%%%%%%%%%%%%%%%%%%%%%%%%%%%%%%%%%%%%%%%%%%%%%%%%%%%%%%
	%  SECCIÓN: Columna Lumbar
	%  Archivo: secciones/columna_lumbar.tex
	%  Sin preámbulo ni \begin{document} — lo maneja main.tex
		%%%%%%%%%%%%%%%%%%%%%%%%%%%%%%%%%%%%%%%%%%%%%%%%%%%%%%%%%%%
		
		\section{Columna Lumbar}
		
		\subsection{Frente AP-PA}
		
		\subsubsection*{1. Identificación}
		
		\textbf{Tipo:} Proyección básica. Forma par con la proyección lateral de columna lumbar.
		
		\textbf{Indicaciones clínicas principales:}
		\begin{itemize}
			\item Fracturas de columna lumbar
			\item Escoliosis y otras curvaturas vertebrales
			\item Procesos neoplásicos vertebrales
		\end{itemize}
		
		\subsubsection*{2. Técnica Base}
		
		\textbf{Receptor de imagen:} 35$\times$43 cm (exploración lumbosacra) o 30$\times$35 cm (solo lumbar). Orientación \textbf{longitudinal (vertical)}.
		
		\textbf{DFRI:} 100 cm.
		
		\textbf{Técnica:} Intermedia, ajustada según el biotipo del paciente. Se utiliza kVp alto para penetrar la zona densa de la columna, permitiendo visualizar los detalles óseos, los espacios interdiscales y las partes blandas. El mAs se ajusta según el biotipo del paciente.
		
		\textbf{Bucky:} Bucky mural o de mesa con grilla.
		
		\subsubsection*{3. Posicionamiento}
		
		\paragraph{Paciente}
		Decúbito supino (posición preferida para la rutina). Plano medio sagital perpendicular al receptor de imagen y centrado en la línea media de la rejilla. Se indica al paciente vaciarse la vejiga urinaria previamente. Se prefiere la posición de bipedestación en pacientes con dolor intenso para reducir las molestias de la exploración. Si el paciente está en decúbito, se lo tracciona desde las axilas (preguntar por prótesis antes de realizar tracción). El paciente debe estar descalzo para evitar diferencias en la altura por calzado.
		
		\begin{figure}[h]
			\centering
			\includegraphics[width=0.7\linewidth]{imagenes/col29}
			\caption{}
			\label{fig:col29}
		\end{figure}
		
		
		\begin{figure}[h]
			\centering
			\includegraphics[width=0.7\linewidth]{imagenes/col30}
			\caption{}
			\label{fig:col30}
		\end{figure}
		
		
		\paragraph{Región de interés}
		Hombros y caderas ajustados en un mismo plano horizontal para evitar rotación. Codos flexionados con manos apoyadas sobre la parte superior del tórax, evitando que los antebrazos entren en el campo de exposición. Se reduce la lordosis lumbar mediante flexión suficiente de caderas y rodillas, colocando la espalda en contacto firme con la mesa. Se puede usar un soporte radiotransparente bajo la parte lateral inferior de la pelvis para reducir la rotación.
		
		\paragraph{Rayo Central}
		Perpendicular al plano del receptor. Para exploración \textbf{solo lumbar}: punto de entrada 4 cm por encima de las crestas ilíacas, incidiendo a nivel de L3. Para exploración \textbf{lumbosacra}: punto de entrada a la altura de las crestas ilíacas, incidiendo a nivel de L4.
		
		\paragraph{Receptor de imagen}
		Centrado en la región lumbar. Debe incluirse desde T12 hasta S1, abarcando una o dos vértebras torácicas inferiores, el sacro-cóccix y los huesos pélvicos. El colimador se ajusta hasta los márgenes laterales de los músculos psoas.
		
		\paragraph{Respiración}
		Suspendida al final de la espiración. La espiración máxima desplaza las vísceras abdominales y reduce la superposición de gas intestinal sobre las estructuras vertebrales, optimizando la nitidez de los cuerpos vertebrales y los espacios discales.
		
		\subsubsection*{4. Anatomía Visible}
		
		Se visualizan los cuerpos vertebrales lumbares (L1--L5), los espacios discales intervertebrales, los espacios interpedunculares, las láminas, y las apófisis espinosas y transversas. Las articulaciones sacroilíacas deben verse equidistantes de la columna vertebral. Las articulaciones intervertebrales deben estar abiertas. En estudios con RI de 35$\times$43 cm se incluyen además una o dos vértebras torácicas inferiores, el sacro, el cóccix y los huesos pélvicos. La sombra del músculo psoas debe ser visible como indicador de buena técnica de penetración.
		
		Para contar las vértebras lumbares se procede de abajo hacia arriba: partiendo de la vértebra ubicada por encima del sacro (L5), se numeran 5, 4, 3, 2 y 1, verificando que L1 sea la que porta costilla.
		
		\textbf{Ventaja de la proyección PA (prono-anterior):} La divergencia del haz de rayos, combinada con la lordosis lumbar, hace que los rayos divergentes pasen tangenciales a los espacios intervertebrales, manteniéndolos bien abiertos. Además reduce la dosis al paciente respecto de la posición AP. En caso de escoliosis detectada en el frente, se apoya el lado convexo contra el RI y el RC entra por el lado cóncavo, aprovechando la divergencia del haz para pasar tangencial a la curvatura.
		
		\subsubsection*{5. Criterios de Evaluación}
		
		\textbf{Posicionamiento correcto:} El área de imagen comprende desde las vértebras torácicas inferiores hasta el sacro. Las apófisis espinosas se visualizan centradas en los cuerpos vertebrales. Las articulaciones intervertebrales aparecen abiertas. La mitad de T12 debe ser visible; de no ser así, se repite la toma.
		
		\textbf{Ausencia de rotación:} Las articulaciones sacroilíacas se encuentran equidistantes de la columna vertebral. Las vértebras son simétricas en ambos lados del plano medio sagital. No debe haber asimetría en las apófisis transversas.
		
		\textbf{Penetración y contraste:} Todos los cuerpos vertebrales lumbares presentan penetración adecuada. La sombra del músculo psoas es visible bilateralmente. El colimado llega hasta los márgenes laterales del psoas, sin incluir partes blandas laterales innecesarias. Ausencia de artefactos producidos por elásticos de ropa interior.
		
		\subsubsection*{6. Puntos Críticos}
		
		\textbf{Rotación del paciente:} Produce asimetría de las articulaciones sacroilíacas y de los pedículos, dificultando la evaluación de los cuerpos vertebrales. Corregir alineando hombros y caderas en el mismo plano horizontal y verificando la perpendicularidad del plano medio sagital al RI.
		
		\textbf{Lordosis lumbar no corregida:} Los espacios intervertebrales aparecen cerrados, especialmente en los niveles L4--L5 y L5--S1, impidiendo su evaluación. Corregir mediante flexión de caderas y rodillas hasta que la espalda contacte firmemente con la mesa.
		
		\textbf{Adaptaciones:} En pacientes con dolor lumbar agudo, utilizar la posición en bipedestación para reducir las molestias. En pacientes con prótesis de cadera u hombro, evitar la tracción de axilas y utilizar otros métodos de estabilización. En pacientes obesos, aumentar el kVp para garantizar la penetración adecuada de la zona densa.
		
		\textbf{Trucos:} Al sospechar o confirmar escoliosis en la imagen, apoyar el lado convexo de la curvatura contra el RI y dirigir el RC desde el lado cóncavo para que los rayos divergentes pasen tangenciales a los espacios intervertebrales de la curva. Verificar siempre la presencia de T12 (con su costilla) para la correcta identificación y conteo vertebral de abajo hacia arriba.
		
		
		\onecolumn
		\clearpage
		
		\begin{figure}
			\centering
			\includegraphics[width=0.5\linewidth]{imagenes/col31}
			\caption{}
			\label{fig:col31}
		\end{figure}
				
		\twocolumn
		\clearpage
		
		%----------------------------------------------------------
		%  Para agregar otro enfoque duplicá el bloque
		%  desde \subsection hasta \textbf{Trucos}
		%----------------------------------------------------------
		
	%%%%%%%%%%%%%%%%%%%%%%%%%%%%%%%%%%%%%%%%%%%%%%%%%%%%%%%%%%%
	%  SECCIÓN: Columna Lumbar
	%  Archivo: secciones/columna_lumbar.tex
	%  Sin preámbulo ni \begin{document} — lo maneja main.tex
		%%%%%%%%%%%%%%%%%%%%%%%%%%%%%%%%%%%%%%%%%%%%%%%%%%%%%%%%%%%
		
		\section{Columna Lumbar}
		
		\subsection{Lateral}
		
		\subsubsection*{1. Identificación}
		
		\textbf{Tipo:} Proyección básica. Forma par con la proyección AP de columna lumbar.
		
		\textbf{Indicaciones clínicas principales:}
		\begin{itemize}
			\item Fractura vertebral y espondilolistesis
			\item Osteoporosis de vértebras lumbares
			\item Procesos neoplásicos
			\item Evaluación de espacios discales y agujeros intervertebrales
		\end{itemize}
		
		\subsubsection*{2. Técnica Base}
		
		\textbf{Receptor de imagen:} 35$\times$43 cm para proyección lumbosacra; 30$\times$35 cm para columna lumbar exclusiva. Orientación \textbf{longitudinal}.
		
		\textbf{DFRI:} 100 cm. Distancia estándar para técnica homogeneizada.
		
		\textbf{Técnica:} Básica con kVp intermedio-alto para penetrar la zona densa de la columna y lograr visualización simultánea de detalles óseos, espacios interdiscales y partes blandas. Para evaluación de litiasis puede optarse por escala corta.
		
		\textbf{Bucky:} Sí. Uso obligatorio de Bucky con grilla por el alto grosor de la región.
		
		\subsubsection*{3. Posicionamiento}
		
		\paragraph{Paciente}
		Decúbito lateral (preferiblemente sobre el lado afectado o derecho para mayor comodidad operativa) o bipedestación lateral. Bata abierta por detrás para dejar la columna expuesta. Plano medio coronal perpendicular al receptor y centrado. El paciente se apoya sobre el lado de interés.
		
		\begin{figure}[h]
			\centering
			\includegraphics[width=0.7\linewidth]{imagenes/col32}
			\caption{}
			\label{fig:col32}
		\end{figure}
		
		
		\paragraph{Región de interés}
		Caderas y rodillas flexionadas cómodamente para reducir la lordosis lumbar. Las rodillas se superponen entre sí para prevenir la rotación pélvica. Se coloca almohadilla bajo la cabeza para mantener la columna horizontal. En pacientes con escoliosis, se apoya el lado convexo contra el receptor y el rayo central ingresa por el lado cóncavo, aprovechando el efecto de abanico para abrir los espacios intervertebrales. La región debe abarcar de T11--T12 hasta S1.
		
		\paragraph{Rayo Central}
		Perpendicular al receptor en condiciones estándar. Punto de entrada: \textbf{a la altura de la cresta ilíaca (L4--L5), 4 traveses de dedo por dentro del reborde cutáneo posterior}. La angulación se modifica según la morfología del paciente: angulación caudal de 5° en hombres y 8° en mujeres (según anchura de pelvis) cuando la columna no está horizontal; mayor angulación caudal ante pelvis más ancha. En columnas rectificadas no se requiere angulación. Para el pasaje L5--S1, se angula con mayor frecuencia; a mayor cadera, mayor angulación necesaria para abrir ese espacio.
		
		\paragraph{Receptor de imagen}
		Centrado a la altura de las crestas ilíacas (L4) para la proyección lumbosacra, o 4 cm por encima de las crestas ilíacas (L3) para la proyección lumbar exclusiva. Los bordes deben incluir desde T11--T12 hasta S1.
		
		\paragraph{Respiración}
		Suspendida al final de una espiración. La espiración completa reduce el diámetro anteroposterior del tronco, disminuye la superposición de estructuras y permite mayor nitidez en la imagen.
		
		\subsubsection*{4. Anatomía Visible}
		
		Se visualizan los cuerpos vertebrales lumbares de perfil, con sus bordes posteriores superpuestos entre sí como criterio de ausencia de rotación. Los espacios intervertebrales deben aparecer abiertos y bien definidos. Los agujeros intervertebrales se muestran abiertos, lo que permite evaluar su morfología. Las apófisis espinosas se proyectan de perfil. Las crestas ilíacas deben aparecer superpuestas cuando el haz no está angulado. Las escotaduras ciáticas mayores también deben superponerse.
		
		\subsubsection*{5. Criterios de Evaluación}
		
		\textbf{Posicionamiento correcto:} Espacios discales intervertebrales y agujeros intervertebrales abiertos. Apófisis espinosas proyectadas de perfil. Región abarcada desde T11--T12 hasta S1.
		
		\textbf{Ausencia de rotación:} Bordes posteriores de cada cuerpo vertebral superpuestos entre sí. Crestas ilíacas superpuestas (en proyección sin angulación). Escotaduras ciáticas mayores superpuestas. Si los cuerpos vertebrales presentan su base superior e inferior desplazadas entre sí (una adelante y otra atrás), indica que el paciente no está completamente de perfil.
		
		\textbf{Penetración y contraste:} Visualización clara y simultánea de estructuras óseas, espacios interdiscales y partes blandas. La técnica de kVp intermedio-alto debe permitir penetrar la zona densa de la columna sin sobreexponer las partes blandas periféricas.
		
		\subsubsection*{6. Puntos Críticos}
		
		\textbf{Rotación corporal:} Produce desalineación de los bordes posteriores vertebrales y de las crestas ilíacas, generando superposición asimétrica que impide la evaluación. Corrección: superponer rodillas, colocar saco de arena entre ellas y verificar que el plano medio coronal quede perpendicular al receptor.
		
		\textbf{Columna no horizontal:} Genera apertura desigual de los espacios intervertebrales. Corrección: colocar apoyo radiotransparente bajo el tórax inferior (método de elección) o aplicar angulación caudal del rayo central (5° en hombres, 8° en mujeres como promedio, ajustando según anchura de pelvis).
		
		\textbf{Adaptaciones:} En pacientes delgados, colocar almohadilla en la cadera declive para horizontalizar la columna. En pacientes con escoliosis, apoyar el lado convexo contra el receptor y dirigir el RC por el lado cóncavo. Para el segmento L5--S1, aumentar la angulación caudal según la anchura de la pelvis del paciente.
		
		\textbf{Trucos:} Colocar el brazo inferior perpendicular al cuerpo para estabilizar al paciente. El saco de arena entre rodillas superpuestas previene la rotación pélvica eficazmente. La rotación residual se detecta observando los cuerpos vertebrales: si la base superior e inferior de un mismo cuerpo no se superponen, el paciente está rotado. En escoliosis, el efecto de abanico al apoyar la convexidad permite abrir los espacios intervertebrales del lado cóncavo.
		
		\onecolumn
		\clearpage
		
		\begin{figure}
			\centering
			\includegraphics[width=0.7\linewidth]{imagenes/col33}
			\caption{}
			\label{fig:col33}
		\end{figure}
		
		\twocolumn
		\clearpage
		
		%----------------------------------------------------------
		%  Para agregar otro enfoque duplicá el bloque
		%  desde \subsection hasta \textbf{Trucos}
		%----------------------------------------------------------
		
	%%%%%%%%%%%%%%%%%%%%%%%%%%%%%%%%%%%%%%%%%%%%%%%%%%%%%%%%%%%
	%  SECCIÓN: Columna Lumbar
	%  Archivo: secciones/columna_lumbar.tex
	%  Sin preámbulo ni \begin{document} — lo maneja main.tex
		%%%%%%%%%%%%%%%%%%%%%%%%%%%%%%%%%%%%%%%%%%%%%%%%%%%%%%%%%%%
		
		\section{Columna Lumbar}
		
		\subsection{Charnela Lumbosacra (L5-S1) - Proyección Lateral}
		
		\subsubsection*{1. Identificación}
		
		\textbf{Tipo:} Complementario.
		
		\textbf{Indicaciones clínicas principales:}
		\begin{itemize}
			\item Espondilolistesis: deslizamiento anterior o posterior de una vértebra respecto a la adyacente, frecuente en el segmento L4-L5 hasta S1.
			\item Espondilosis: degeneración discovertebal que debilita el hueso, afecta principalmente L5-S1.
			\item Discopatías y estenosis del canal en la charnela lumbosacra.
		\end{itemize}
		
		\subsubsection*{2. Técnica Base}
		
		\textbf{Receptor de imagen:} 18$\times$24 cm. Orientación \textbf{longitudinal} (vertical), adaptada al eje mayor de la región lumbosacra.
		
		\textbf{DFRI:} 100 cm. Distancia estándar para columna lumbar.
		
		\textbf{Técnica:} Intermedia en pacientes estándar (kVp alto para penetrar la zona densa de la columna y visualizar detalles óseos, espacios interdiscales y partes blandas). Básica en pacientes con mayor grosor corporal. Debe compensarse la colimación reducida aplicando un incremento del 50--75\% en los factores, calculado mediante regla de tres.
		
		\textbf{Bucky:} Sí, con rejilla antidifusora (Bucky de mesa).
		
		\subsubsection*{3. Posicionamiento}
		
		\paragraph{Paciente}
		Decúbito lateral sobre la mesa radiológica. En caso de inestabilidad o sospecha de trauma, puede realizarse de pie. El plano medio sagital (PMS) de la cabeza debe ubicarse en el mismo plano que el eje longitudinal de la columna, utilizando una almohada de apoyo.
		
		\begin{figure}[h]
			\centering
			\includegraphics[width=0.7\linewidth]{imagenes/col34}
			\caption{}
			\label{fig:col34}
		\end{figure}
	
		\paragraph{Región de interés}
		El plano medio coronal (PMC) del cuerpo debe quedar perpendicular al receptor de imagen. Se superponen las rodillas con un apoyo entre ellas para evitar la rotación pélvica. Si es posible, las caderas se extienden completamente. Se puede colocar un soporte radiotransparente bajo la parte inferior del tórax para que el eje longitudinal de la columna quede en posición horizontal verdadera, lo que es el método de elección. El brazo inferior se flexiona en el codo y se posiciona perpendicular al cuerpo.
		
		\paragraph{Rayo Central}
		Perpendicular al receptor cuando la columna está correctamente horizontalizada. Punto de entrada: \textbf{5 cm posterior a la EIAS elevada y 4 cm inferior a la cresta ilíaca}. Para localizar el PMC se apoya una mano sobre el reborde cutáneo posterior del paciente y se avanzan 4 dedos hacia adelante; para la altura, se palpa la cresta ilíaca y se descienden 2 dedos.
		
			
		
		\begin{figure}[h]
			\centering
			\includegraphics[width=0.7\linewidth]{imagenes/col35}
			\caption{}
			\label{fig:col35}
		\end{figure}
		
		
		
		Cuando la columna no puede horizontalizarse, se angula el RC en sentido \textbf{caudal 5° en hombres} y \textbf{8° en mujeres}, ingresando por la EIAS, para mejorar la apertura del espacio intervertebral L5-S1.
		
		\textit{Técnica alternativa de Francis:} localizar ambas crestas ilíacas, trazar una línea imaginaria entre ellas (plano interilíaco) y dirigir el RC paralelo a dicha línea.
		
		\begin{figure}[h]
			\centering
			\includegraphics[width=0.7\linewidth]{imagenes/col37}
			\caption{}
			\label{fig:col37}
		\end{figure}
		
		
		\textit{Referencia angular adicional:} el espacio L5-S1 presenta una orientación oblicua de 10--15° condicionada por la lordosis lumbar y la inclinación sacra (35--50° respecto a la horizontal). Esto implica que, para obtener incidencia perpendicular al espacio intervertebral sin horizontalización, serían necesarios 30--35° en sentido craneal o 10--15° en sentido caudal según la variante anatómica individual.
		
		\paragraph{Receptor de imagen}
		Centrado en la articulación L5-S1. Colimación estrecha sobre la charnela lumbosacra. El área de exposición debe incluir los cuerpos vertebrales de L5 y S1 completos, y la parte superior del sacro.
		
		\paragraph{Respiración}
		Suspendida durante la exposición para eliminar microdesplazamientos de las estructuras óseas adyacentes.
		
		\subsubsection*{4. Anatomía Visible}
		
		Se visualiza la unión lumbosacra (L5-S1) con el espacio intervertebral completamente abierto, la vértebra L5 íntegra y la parte superior del sacro (S1). Las apófisis espinosas se proyectan de perfil. Las crestas ilíacas deben aparecer casi superpuestas entre sí cuando el RC no está angulado. La charnela lumbosacra ocupa el centro del área de exposición.
		
		\subsubsection*{5. Criterios de Evaluación}
		
		\textbf{Posicionamiento correcto:} articulación intervertebral lumbosacra L5-S1 completamente abierta; toda la vértebra L5 y la parte superior del sacro incluidas dentro del área colimada; charnela lumbosacra centrada en el área de exposición.
		
		\textbf{Ausencia de rotación:} crestas ilíacas casi superpuestas entre sí (cuando el RC no está angulado); apófisis espinosas proyectadas de perfil sin oblicuidad.
		
		\textbf{Penetración y contraste:} adecuada visualización de las estructuras óseas densas de la charnela; contraste suficiente para diferenciar los espacios discales y las partes blandas perivertebral.
		
		\subsubsection*{6. Puntos Críticos}
		
		\textbf{Espacio articular L5-S1 cerrado:} consecuencia de no aplicar la angulación caudal cuando la columna no está en posición horizontal verdadera. Corregir horizontalizando la columna con soporte radiotransparente bajo el tórax inferior, o angulando el RC 5° (hombres) u 8° (mujeres) en sentido caudal.
		
		\textbf{Rotación pélvica:} impide la superposición de las crestas ilíacas y genera proyección oblicua de las apófisis espinosas. Corregir verificando que las rodillas estén exactamente superpuestas con apoyo intermedio antes de exponer.
		
		\textbf{Adaptaciones:} en pacientes con trauma o inestabilidad, mantener la posición lateral sin movilizaciones excesivas; puede realizarse de pie si es necesario. Cuando no se pueda horizontalizar la columna, aplicar la técnica alternativa de Francis. En pacientes con cintura grande puede requerirse angulación cefálica del RC.
		
		\textbf{Trucos:} la EIAS elevada es fácilmente palpable con el paciente en decúbito lateral y constituye el punto de referencia estandarizado para centrar en L5-S1. El soporte radiotransparente bajo el tórax es el método de elección para horizontalizar la columna. Compensar la colimación estrecha incrementando los factores técnicos un 50--75\% (calcular mediante regla de tres). Verificar superposición de rodillas antes de la exposición.
		
		\onecolumn
		\clearpage
		
		\begin{figure}
			\centering
			\includegraphics[width=0.7\linewidth]{imagenes/col36}
			\caption{}
			\label{fig:col36}
		\end{figure}
		
		
		\twocolumn
		\clearpage
		
		%----------------------------------------------------------
		%  Para agregar otro enfoque duplicá el bloque
		%  desde \subsection hasta \textbf{Trucos}
		%----------------------------------------------------------
		
	%%%%%%%%%%%%%%%%%%%%%%%%%%%%%%%%%%%%%%%%%%%%%%%%%%%%%%%%%%%
	%  SECCIÓN: Columna Lumbar
	%  Archivo: secciones/columna_lumbar_oblicua.tex
	%  Sin preámbulo ni \begin{document} — lo maneja main.tex
		%%%%%%%%%%%%%%%%%%%%%%%%%%%%%%%%%%%%%%%%%%%%%%%%%%%%%%%%%%%
		
		\section{Columna Lumbar}
		
		\subsection{AP Oblicua (OPD / OPI)}
		
		\subsubsection*{1. Identificación}
		
		\textbf{Tipo:} Proyección complementaria. Se realizan ambas oblicuas en conjunto de forma comparativa bilateral.
		
		\textbf{Indicaciones clínicas principales:}
		\begin{itemize}
			\item Evaluación de articulaciones interapofisarias y pars interarticularis (espondilólisis, espondilolistesis)
			\item Dolor crónico de columna; también indicada como referencia anatómica previa a infiltración anestésica interapofisaria bajo arco en C
			\item Fracturas de pars interarticularis no visibles en proyecciones estándar
		\end{itemize}
		
		\subsubsection*{2. Técnica Base}
		
		\textbf{Receptor de imagen:} 35$\times$43 cm (o 30$\times$35 cm). Orientación \textbf{longitudinal}. Para estudio focalizado en la quinta articulación interapofisaria se reduce a 18$\times$24 cm.
		
		\textbf{DFRI:} 100 cm. Distancia estándar para columna lumbar que equilibra nitidez geométrica y magnificación.
		
		\textbf{Técnica:} Intermedia para paciente estándar; básica para paciente de mayor grosor corporal. Se emplea kVp alto para garantizar la penetración de la zona densa de la columna y obtener visibilidad adecuada de detalles óseos, espacios interdiscales y partes blandas. \textbf{kVp --- mAs:} \textit{[dato no suministrado]}.
		
		\textbf{Bucky:} Con Bucky de mesa.
		
		\subsubsection*{3. Posicionamiento}
		
		\paragraph{Paciente}
		Decúbito supino sobre la mesa radiográfica, o en bipedestación según capacidad del paciente. El decúbito es preferible: estabiliza la posición evitando que la pelvis rote independientemente del tronco, y en pacientes de gran tamaño corporal reduce el espesor efectivo, disminuyendo la radiación dispersa y mejorando la calidad de imagen.
		
		\begin{figure}[h]
			\centering
			\includegraphics[width=0.7\linewidth]{imagenes/col38OPD}
			\caption{OPD}
			\label{fig:col38-opd}
		\end{figure}
		
		
		\paragraph{Región de interés}
		Partiendo del decúbito supino con el plano medio sagital perpendicular al receptor, se rota el cuerpo hacia el lado a estudiar aproximadamente \textbf{45°}. La \textbf{OPD} (rotación hacia la derecha) demuestra las articulaciones interapofisarias derechas; la \textbf{OPI} (rotación hacia la izquierda) demuestra las izquierdas. En ambos casos se visualizan las articulaciones del lado más próximo al receptor. Para las articulaciones de L5--S1 puede requerirse hasta \textbf{60°} de oblicuidad. Las articulaciones lumbares forman un ángulo de 30--60° con el plano sagital, por lo que la angulación varía según el nivel: menor angulación para articulaciones más superiores, mayor para las más inferiores. El eje longitudinal del paciente se mantiene paralelo al eje de la mesa. Se pueden colocar soportes bajo el hombro elevado, la cadera y la rodilla para evitar movimientos involuntarios.
		
		\paragraph{Rayo Central}
		Perpendicular al plano del receptor. Punto de entrada estándar: \textbf{5 cm medial a la EIAS elevada y 4 cm por encima de la cresta ilíaca} (nivel aproximado de L3). Para énfasis en la quinta articulación interapofisaria: \textbf{5 cm medial a la EIAS elevada y en el punto medio entre la cresta ilíaca y la EIAS}.
		
		\paragraph{Receptor de imagen}
		Centrado en la región lumbar media. El campo debe abarcar desde las vértebras torácicas inferiores hasta el sacro. La columna vertebral debe quedar paralela a la superficie del receptor para que los espacios intervertebrales T12--L1 y L1--L2 permanezcan abiertos.
		
		\paragraph{Respiración}
		Suspendida al final de una espiración completa. La espiración reduce el diámetro abdominal, disminuye la superposición de estructuras y minimiza el movimiento durante la exposición.
		
		\subsubsection*{4. Anatomía Visible}
		
		Se visualizan las articulaciones interapofisarias del lado más próximo al receptor, abiertas y uniformemente visibles a través de los cuerpos vertebrales. Se identifican las apófisis articulares superiores e inferiores y la pars interarticularis, estructura clave para detectar fracturas en espondilólisis y espondilolistesis.
		
		El criterio anatómico central es la imagen del \textbf{perro Scottie} (Proyección de La Chapelle): cuando el posicionamiento es correcto, cada vértebra lumbar adquiere la silueta de un perro escocés. Sus componentes son: el \textbf{ojo} corresponde al pedículo; la \textbf{oreja} a la apófisis articular superior (que articula con la apófisis articular inferior de la vértebra suprayacente, formando la ``pata''); el \textbf{hocico} a la apófisis transversa; el \textbf{cuello} a la pars interarticularis; y el \textbf{cuerpo} a la lámina. Una fractura de la pars interarticularis se evidencia como una línea radiolúcida en el cuello del perro, visible únicamente en proyección oblicua.
		
		Se aprecian además los cuerpos vertebrales, pedículos, espacios interdiscales y partes blandas paravertebrales.
		
		\begin{figure}[h]
			\centering
			\includegraphics[width=0.7\linewidth]{imagenes/col39}
			\caption{}
			\label{fig:col39}
		\end{figure}
		
		
		\subsubsection*{5. Criterios de Evaluación}
		
		\textbf{Posicionamiento correcto:} La figura del perro Scottie se forma con nitidez en cada vértebra lumbar. Las articulaciones interapofisarias del lado estudiado se encuentran abiertas y visibles de forma uniforme. El pedículo se proyecta centrado sobre el cuerpo vertebral. Los espacios intervertebrales T12--L1 y L1--L2 están abiertos, confirmando que la columna es paralela al receptor.
		
		\textbf{Ausencia de rotación:} La oblicuidad es correcta cuando el pedículo se ubica sobre el cuerpo vertebral sin desplazamiento anterior ni posterior apreciable. La articulación interapofisaria aparece abierta y bien definida.
		
		\textbf{Penetración y contraste:} El kVp alto produce escala larga de grises que permite visualizar simultáneamente los detalles trabeculares óseos, los espacios interdiscales y las partes blandas adyacentes en una sola exposición.
		
		\subsubsection*{6. Puntos Críticos}
		
		\textbf{Rotación insuficiente:} El pedículo (ojo del perro) se proyecta anterior al cuerpo vertebral y la articulación interapofisaria no se abre. Corrección: aumentar la oblicuidad del paciente hacia el lado estudiado.
		
		\textbf{Rotación excesiva:} El pedículo se proyecta posterior al cuerpo vertebral. Corrección: reducir la oblicuidad del paciente.
		
		\textbf{Adaptaciones:} Para L5--S1 se incrementa la oblicuidad hasta 60° y, si el estudio es exclusivo de esa articulación, se utiliza RI de 18$\times$24 cm con RC en el punto medio entre la cresta ilíaca y la EIAS elevada. En pacientes de gran tamaño corporal, el decúbito es preferible al bipedestación por la estabilidad posicional y la reducción de la dispersión.
		
		\textbf{Trucos:} La identificación de la marca lateral sigue esta lógica: en una OPD los cuerpos vertebrales (estructura más anterior) se orientan hacia el lado derecho de la imagen y las apófisis espinosas hacia el lado opuesto; en la OPI, a la inversa. La comparación sistemática OPD--OPI es indispensable para diferenciar espondilólisis unilateral de espondilolistesis bilateral. En el contexto de infiltración interapofisaria, ambas oblicuas se solicitan en bloque y sirven como referencia anatómica previa al procedimiento bajo arco en C.
		
		
		\onecolumn
		\clearpage
		
		\begin{figure}
			\centering
			\includegraphics[width=0.7\linewidth]{imagenes/col40}
			\caption{}
			\label{fig:col40}
		\end{figure}
		
		
		\twocolumn
		\clearpage
		
		%----------------------------------------------------------
		%  Para agregar otro enfoque duplicá el bloque
		%  desde \subsection hasta \textbf{Trucos}
		%----------------------------------------------------------
		
	
	%%%%%%%%%%%%%%%%%%%%%%%%%%%%%%%%%%%%%%%%%%%%%%%%%%%%%%%%%%%
	%  SECCIÓN: Columna Lumbar
	%  Archivo: secciones/columna_lumbar_pa_oblicua.tex
	%  Sin preámbulo ni \begin{document} — lo maneja main.tex
		%%%%%%%%%%%%%%%%%%%%%%%%%%%%%%%%%%%%%%%%%%%%%%%%%%%%%%%%%%%
		
		\section{Columna Lumbar}
		
		\subsection{PA Oblicua (OAD / OAI)}
		
		\subsubsection*{1. Identificación}
		
		\textbf{Tipo:} Proyección complementaria. OAD y OAI son el mismo enfoque ejecutado hacia cada lado; se realizan en par para estudiar ambas columnas interapofisarias.
		
		\textbf{Indicaciones clínicas principales:}
		\begin{itemize}
			\item Defectos en la pars interarticularis (espondilólisis)
			\item Espondilolistesis
			\item Evaluación de las articulaciones interapofisarias del lado más alejado del RI
		\end{itemize}
		
		\subsubsection*{2. Técnica Base}
		
		\textbf{Receptor de imagen:} 35$\times$43 cm. Orientación \textbf{longitudinal} en paciente estándar. Para estudio exclusivo de la última articulación interapofisaria: 18$\times$24 cm.
		
		\textbf{DFRI:} 100 cm. Distancia estándar para radiografía de columna con Bucky de mesa o mural.
		
		\textbf{Técnica:} Intermedia en paciente estándar; básica (mayor carga) en paciente de mayor espesor.
		
		\textbf{Bucky:} Bucky con grilla. Uso obligatorio dado el espesor de la región lumbar.
		
		\subsubsection*{3. Posicionamiento}
		
		\paragraph{Paciente}
		Decúbito prono (posición preferida por mayor facilidad de inmovilización y menor porcentaje de repeticiones respecto a la oblicua AP). Alternativa: posición incorporada. Desde el decúbito prono, se rota al paciente hacia una posición de semiprono, apoyando el cuerpo sobre el antebrazo y la rodilla flexionados del lado que quedará elevado.
		
		\begin{figure}[h]
			\centering
			\includegraphics[width=0.7\linewidth]{imagenes/col41}
			\caption{OAI}
			\label{fig:col41}
		\end{figure}
		
		
		\paragraph{Región de interés}
		Rotación corporal oblicua de 45°; puede requerirse hasta 60° para demostrar correctamente las articulaciones interapofisarias de L5--S1 (charnela lumbosacra). La pierna del lado elevado se flexiona formando un ``4'' con la pierna contralateral extendida. Centrar L3 en la línea media de la rejilla. Para demostrar la charnela lumbosacra, centrar el RI sobre L5. La columna vertebral debe quedar paralela a la superficie de la mesa para mantener abiertos los espacios intervertebrales T12--L1 y L1--L2.
		
		\paragraph{Rayo Central}
		Perpendicular al plano del receptor. Punto de entrada: \textbf{L3, entre 2,5 y 3,8 cm por encima de la cresta ilíaca}, aproximadamente 5 cm lateral a las apófisis espinosas palpables. El RC entra siempre por el \textbf{lado elevado}.
		
		\paragraph{Receptor de imagen}
		Centrado en L3. La región a incluir abarca desde las vértebras torácicas inferiores (T12--L1) hasta el sacro. En la variante de 18$\times$24 cm, centrar sobre la última articulación interapofisaria de interés.
		
		\paragraph{Respiración}
		Suspendida al final de la espiración. La espiración completa reduce el espesor abdominal y la superposición de gas intestinal, mejorando el contraste de los elementos óseos posteriores.
		
		\subsubsection*{4. Anatomía Visible}
		
		Se demuestran las articulaciones interapofisarias del lado \textbf{más alejado del RI} (lado contrario al apoyo). En OAD el paciente apoya el lado derecho y se visualizan las articulaciones izquierdas; en OAI apoya el lado izquierdo y se visualizan las articulaciones derechas. Se incluyen las apófisis articulares superiores e inferiores, la pars interarticularis y la articulación T12--L1 en el RI de 35$\times$43 cm. La quinta articulación lumbosacra suele quedar bien visible. Los espacios intervertebrales T12--L1 y L1--L2 deben estar abiertos.
	
		Imagen del ``perro de Scottie'': figura radiológica característica de la oblicua lumbar cuyas partes corresponden a estructuras anatómicas específicas:
		\begin{itemize}
			\item \textbf{Ojo:} pedículo
			\item \textbf{Oreja:} apófisis articular superior
			\item \textbf{Hocico:} apófisis transversa
			\item \textbf{Cuello:} pars interarticularis
			\item \textbf{Cuerpo:} lámina
			\item \textbf{Patas delanteras:} apófisis articular inferior
		\end{itemize}
		
			
		\begin{figure}[h]
			\centering
			\includegraphics[width=0.7\linewidth]{imagenes/col39}
			\caption{}
			\label{fig:col98988}
		\end{figure}
		
		
		
		\textbf{Marca y orientación del cuerpo vertebral:}
		\begin{itemize}
			\item OAD: el cuerpo vertebral mira hacia la \textbf{izquierda} (lado contrario a la marca).
			\item OAI: el cuerpo vertebral mira hacia la \textbf{derecha} (mismo lado que la marca).
		\end{itemize}
		
		\subsubsection*{5. Criterios de Evaluación}
		
		\textbf{Posicionamiento correcto:} Articulaciones interapofisarias del lado más alejado del RI claramente abiertas y visibles. Columna vertebral paralela a la superficie de la mesa, con apertura de los espacios intervertebrales T12--L1 y L1--L2. Área cubierta desde las vértebras torácicas inferiores hasta el sacro.
		
		\textbf{Rotación insuficiente o excesiva:} Si la articulación interapofisaria no se visualiza bien y el pedículo aparece claramente anterior al cuerpo vertebral, la rotación es insuficiente (aumentar la oblicuidad). Si el pedículo aparece claramente posterior al cuerpo vertebral, la rotación es excesiva (reducir la oblicuidad).
		
		\textbf{Penetración y contraste:} Técnica adecuada cuando se distinguen con claridad las trabéculas óseas y los espacios articulares. Ajustar a técnica básica en pacientes de mayor espesor.
		
		\subsubsection*{6. Puntos Críticos}
		
		\textbf{Rotación insuficiente:} Las articulaciones interapofisarias no se abren y el pedículo se proyecta anterior al cuerpo vertebral. Corrección: aumentar la oblicuidad, hasta 60° si es necesario para L5--S1.
		
		\textbf{Aumento de la DORI:} La posición de semiprono incrementa la distancia objeto--receptor de imagen respecto a la oblicua AP, lo que puede reducir la resolución espacial. Es la principal desventaja de esta proyección frente a la AP oblicua.
		
		\textbf{Adaptaciones:} En paciente con limitaciones de movilidad o trauma, puede realizarse en posición incorporada. Verificar que se mantenga la oblicuidad de 45° con apoyo de cuñas radiolúcidas si es necesario.
		
		\textbf{Trucos:} Pedir al paciente que forme el número ``4'' con la pierna elevada estabiliza la rotación y facilita reproducir el ángulo. Regla práctica: en la PA oblicua lumbar siempre se visualizan las articulaciones del lado \textit{contrario} al apoyo, lo opuesto a las oblicuas torácicas.
		
		
		\onecolumn
		\clearpage
		
		\begin{figure}
			\centering
			\includegraphics[width=0.7\linewidth]{imagenes/col42}
			\caption{OAI/OPD}
			\label{fig:col42}
		\end{figure}
		
		
		\twocolumn
		\clearpage
		
		%----------------------------------------------------------
		%  Para agregar otro enfoque duplicá el bloque
		%  desde \subsection hasta \textbf{Trucos}
		%----------------------------------------------------------
	
	
	%%%%%%%%%%%%%%%%%%%%%%%%%%%%%%%%%%%%%%%%%%%%%%%%%%%%%%%%%%%
	%  SECCIÓN: [Nombre de la Sección]
	%  Archivo: secciones/[nombre_archivo].tex
	%  Sin preámbulo ni \begin{document} — lo maneja main.tex
		%%%%%%%%%%%%%%%%%%%%%%%%%%%%%%%%%%%%%%%%%%%%%%%%%%%%%%%%%%%
		
		\section{Columna Lumbar}
		
		\subsection{Proyecciones Funcionales Lumbares - Hiperflexión / Hiperextensión}
		
		\subsubsection*{1. Identificación}
		
		\textbf{Tipo:} Complementaria. Forma par entre sí: se obtienen dos imágenes laterales, una en hiperflexión y otra en hiperextensión.
		
		\textbf{Indicaciones clínicas principales:}
		\begin{itemize}
			\item Valoración de la movilidad en sitio de artrodesis vertebral.
			\item Inestabilidad vertebral y subluxaciones.
			\item Localización de hernia discal (evidenciada por limitación del movimiento en el punto de la lesión).
		\end{itemize}
		
		\subsubsection*{2. Técnica Base}
		
		\textbf{Receptor de imagen:} 35$\times$43 cm. Orientación \textbf{longitudinal}, colocado en vertical. Borde inferior del RI 3--5 cm por debajo de la cresta ilíaca.
		
		\textbf{DFRI:} 100 cm. Distancia estándar para la región en posición de perfil.
		
		\textbf{Técnica:} Contrastada (alto contraste) con kV moderado. El alto contraste permite definir con precisión los bordes óseos para detectar posibles desplazamientos vertebrales; el kV moderado garantiza la penetración adecuada de la región en posición lateral durante los movimientos de flexo-extensión.
		
		\textbf{Bucky:} Sí, con rejilla.
		
		\subsubsection*{3. Posicionamiento}
		
		\paragraph{Paciente}
		De pie, descalzo y en carga con peso distribuido equitativamente entre ambos pies. Al ser una proyección funcional debe realizarse de pie. Como alternativa puede realizarse en decúbito lateral. La pelvis permanece inmóvil durante toda la exploración.
		
		\begin{figure}[h]
			\centering
			\includegraphics[width=0.7\linewidth]{imagenes/col43}
			\caption{Flexión}
			\label{fig:col43}
		\end{figure}
		
		\begin{figure}
			\centering
			\includegraphics[width=0.7\linewidth]{imagenes/col44}
			\caption{Extensión}
			\label{fig:col44}
		\end{figure}
		
		
		\paragraph{Región de interés}
		Plano medio sagital paralelo al receptor. Se incluye desde T12 hasta la parte superior del cóccix, con bordes posterior y anterior sin cortar las apófisis espinosas.
		
		\paragraph{Rayo Central}
		Perpendicular al plano del receptor. Punto de entrada: \textbf{4 cm por encima de las crestas ilíacas}.
		
		\paragraph{Receptor de imagen}
		Vertical (longitudinal). Centrado a nivel del RC. Borde inferior 3--5 cm por debajo de la cresta ilíaca. Colimación a los 4 lados del área de interés.
		
		\paragraph{Respiración}
		Suspendida en espiración. Favorece la relajación muscular y facilita los movimientos de flexión y extensión. Se realizan dos exposiciones separadas, una por posición.
		
		\subsubsection*{4. Anatomía Visible}
		
		Cuerpos vertebrales de T12, L1--L5, sacro y parte superior del cóccix. Discos intervertebrales, espacios discales y apófisis espinosas de perfil. Articulaciones interapofisiarias.
		
		En \textbf{hiperflexión}: mayor separación entre apófisis espinosas y menor espacio entre cuerpos vertebrales. En \textbf{hiperextensión}: menor espacio entre apófisis espinosas y mayor espacio entre cuerpos vertebrales.
		
		\subsubsection*{5. Criterios de Evaluación}
		
		\textbf{Posicionamiento correcto:} En hiperflexión, el paciente se inclina hacia adelante sin doblar rodillas ni desplazar la pelvis hacia atrás; en decúbito lateral, muslos elevados al máximo. En hiperextensión, el paciente se inclina hacia atrás como intentando mirar el techo, con pies firmes en el suelo. Los marcadores identificadores (hiperflexión / hiperextensión) deben estar correctamente utilizados en cada proyección.
		
		\textbf{Ausencia de rotación:} Márgenes posteriores de los cuerpos vertebrales superpuestos, sin rotación de la columna vertebral.
		
		\textbf{Penetración y contraste:} Bordes óseos de los cuerpos vertebrales y apófisis espinosas bien definidos con alto contraste. Espacios discales y apófisis espinosas claramente diferenciables en ambas posiciones, permitiendo valorar la movilidad en el sitio de artrodesis o restricción por lesión.
		
		\subsubsection*{6. Puntos Críticos}
		
		\textbf{Rotación vertebral:} Produce superposición asimétrica de los márgenes vertebrales, impidiendo la correcta valoración de la movilidad. Corregir alineando el plano medio sagital paralelo al RI y verificando que pelvis y hombros queden en el mismo plano.
		
		\textbf{Insuficiente amplitud de movimiento:} Imágenes con escasa diferencia entre ambas posiciones no permiten valorar la movilidad articular. Guiar al paciente verbalmente con referencias concretas (``doble la espalda como si quisiera tocarse los pies'' / ``inclínese hacia atrás como si quisiera mirar el techo'') sin forzar el movimiento.
		
		\textbf{Adaptaciones:} Si el paciente no tolera la bipedestación, puede realizarse en decúbito lateral; en esa posición se le solicita llevar los muslos hacia arriba (hiperflexión) o extender el torso y miembros hacia atrás (hiperextensión). Para la hiperextensión de pie puede utilizarse un soporte de suero, barras o barandas para que el paciente se afirme. Todos los movimientos deben ser cómodos y dentro del rango tolerable.
		
		\textbf{Trucos:} Preferir siempre la posición erguida por ser más fisiológica y reproducible. Utilizar marcadores identificadores en cada imagen para distinguir hiperflexión de hiperextensión. Valorar comparativamente los cambios en los espacios discales y en las apófisis espinosas entre ambas proyecciones para determinar si existe movilidad en el sitio de fusión o restricción por lesión.
		
		
		\clearpage
		\onecolumn
		
		\begin{figure}
			\centering
			\includegraphics[width=0.7\linewidth]{imagenes/col45}
			\caption{}
			\label{fig:col45}
		\end{figure}
		
		
		\twocolumn
		\clearpage
		%----------------------------------------------------------
		%  Para agregar otro enfoque duplicá el bloque
		%  desde \subsection hasta \textbf{Trucos}
		%----------------------------------------------------------
		
		
	%%%%%%%%%%%%%%%%%%%%%%%%%%%%%%%%%%%%%%%%%%%%%%%%%%%%%%%%%%%
	%  SECCIÓN: Articulaciones Sacroilíacas
	%  Archivo: secciones/sacroiliacas_ap_pa_axial.tex
	%  Sin preámbulo ni \begin{document} — lo maneja main.tex
		%%%%%%%%%%%%%%%%%%%%%%%%%%%%%%%%%%%%%%%%%%%%%%%%%%%%%%%%%%%
		
		\section{Columna Sacra}
		
		\subsection{Proyección AP/PA Axial - Método de Ferguson}
		
		\subsubsection*{1. Identificación}
		
		\textbf{Tipo:} Proyección básica. Las variantes AP (supino) y PA (prono) son intercambiables según la condición del paciente; la PA en prono es la recomendada por Meese para las articulaciones sacroilíacas.
		
		\textbf{Indicaciones clínicas principales:}
		\begin{itemize}
			\item Fracturas de la articulación sacroilíaca
			\item Luxaciones y subluxaciones sacroilíacas
			\item Evaluación de la charnela lumbosacra (pasaje L5--S1)
		\end{itemize}
		
		\subsubsection*{2. Técnica Base}
		
		\textbf{Receptor de imagen:} 18$\times$24 cm o 24$\times$30 cm. Orientación \textbf{longitudinal} respecto al eje mayor de la región.
		
		\textbf{DFRI:} 100 cm.
		
		\textbf{Técnica:} Intermedia en paciente estándar --- Básica en paciente con gran espesor --- Homogenizada en niños o pacientes muy delgados.
		
		\textbf{Bucky:} Bucky de mesa con grilla.
		
		\subsubsection*{3. Posicionamiento}
		
		\paragraph{Paciente}
		\textbf{AP axial:} decúbito supino, plano medio sagital centrado en la rejilla.
		\textbf{PA axial:} decúbito prono (recomendado por Meese), plano medio sagital centrado en la rejilla. En prono, la oblicuidad natural de las articulaciones sacroilíacas las ubica más paralelas al receptor, y la divergencia del haz mejora su visualización sin necesidad de mayor angulación.
		
		\begin{figure}[h]
			\centering
			\includegraphics[width=0.7\linewidth]{imagenes/col46}
			\caption{}
			\label{fig:col46}
		\end{figure}
		
		\begin{figure}[h]
			\centering
			\includegraphics[width=0.7\linewidth]{imagenes/col48}
			\caption{}
			\label{fig:col48}
		\end{figure}
				
		
		\paragraph{Región de interés}
		En ambas variantes se extienden los miembros inferiores o se abducen los muslos, ajustándolos en posición vertical. Debe verificarse que la pelvis no esté rotada. La combinación de la oblicuidad articular y la divergencia del haz permite que el rayo incida de forma tangencial sobre las articulaciones sacroilíacas, mejorando su visualización. Se protegen las gónadas.
		
		\paragraph{Rayo Central}
		\textbf{AP axial:} angulado en dirección \textbf{cefálica}. En hombres: \textbf{30°}. En mujeres: \textbf{35°} (mayor lordosis y ángulo lumbosacro). Ferguson recomendó originalmente 45°; el rango de 30--35° es actualmente el estándar. Dirigido a través de la articulación lumbosacra. Punto de entrada: \textbf{4 cm por encima de la sínfisis del pubis, en el plano medio sagital}.
		
		\textbf{PA axial:} angulado en dirección \textbf{caudal}, con una angulación promedio de \textbf{35°}. Dirigido a través de la articulación lumbosacra hacia el punto medio del RI. Punto de entrada: \textbf{apófisis espinosa de L4}.
		
		La justificación de la angulación en ambos casos es la misma: las curvaturas de la columna impiden que el rayo pase tangencial al espacio articular sin angular; de ahí el calificativo \textit{axial}.
		
		\paragraph{Receptor de imagen}
		Centrado sobre el rayo central en ambas variantes. Los márgenes laterales deben incluir ambas articulaciones sacroilíacas sin cortarlas. La colimación se restringe al área de interés asegurando la inclusión de la charnela lumbosacra y el sacro completo.
		
		\paragraph{Respiración}
		Suspendida en el momento de la exposición. Evita el movimiento de estructuras abdominopélvicas que generarían borrosidad.
		
		\subsubsection*{4. Anatomía Visible}
		
		La imagen muestra la charnela lumbosacra (L5--S1) y el sacro. Se visualizan simétricamente ambas articulaciones sacroilíacas, libres de superposición. La angulación abre el espacio intervertebral entre L5 y S1, que de otra forma quedaría obliterado por las curvaturas de la columna. En la variante PA (Meese), el posicionamiento en prono acerca las articulaciones al plano paralelo al receptor, optimizando aún más su representación.
		
		\subsubsection*{5. Criterios de Evaluación}
		
		\textbf{Posicionamiento correcto:} espacio intervertebral L5--S1 abierto y visible; charnela lumbosacra y sacro incluidos en su totalidad.
		
		\textbf{Ausencia de rotación:} imagen simétrica de ambas articulaciones sacroilíacas, sin superposiciones que las oculten.
		
		\textbf{Penetración y contraste:} ambas articulaciones sacroilíacas adecuadamente penetradas, con densidad que permita distinguir el espacio articular.
		
		\subsubsection*{6. Puntos Críticos}
		
		\textbf{Angulación insuficiente:} el espacio intervertebral L5--S1 no se abre y las articulaciones sacroilíacas quedan superpuestas. Corregir ajustando los grados según el sexo del paciente y valorando visualmente el contorno de la zona lumbar baja para estimar la lordosis real.
		
		\textbf{Rotación pélvica:} asimetría de las articulaciones sacroilíacas en la imagen. Verificar que las EIAS equidisten del plano medio sagital antes de exponer.
		
		\textbf{Adaptaciones:} en pacientes con dolor o trauma agudo que no toleran el decúbito prono, utilizar la variante AP axial. Valorar la angulación según la lordosis real del paciente, aumentándola o reduciéndola respecto al promedio de 30--35°.
		
		\textbf{Trucos:} valorar el contorno de la parte inferior de la espalda para estimar una acentuación o disminución no habitual del ángulo lumbosacro y variar la angulación del RC en consecuencia. En la variante PA, el RC entra por la línea media aproximadamente 5 cm distal a la apófisis espinosa de L5, centrado a la altura de las EIAS.
		
		\clearpage
		\onecolumn
		
		\begin{figure}
			\centering
			\includegraphics[width=0.7\linewidth]{imagenes/col47}
			\caption{AP}
			\label{fig:col47}
		\end{figure}
		
		\begin{figure}
			\centering
			\includegraphics[width=0.7\linewidth]{imagenes/col49}
			\caption{PA}
			\label{fig:col49}
		\end{figure}
		
		
		\twocolumn
		\clearpage
		
		%----------------------------------------------------------
		%  Para agregar otro enfoque duplicá el bloque
		%  desde \subsection hasta \textbf{Trucos}
		%----------------------------------------------------------
	
	
	%%%%%%%%%%%%%%%%%%%%%%%%%%%%%%%%%%%%%%%%%%%%%%%%%%%%%%%%%%%
	%  SECCIÓN: Articulaciones Sacroilíacas
	%  Archivo: secciones/sacroiliacas_oblicua.tex
	%  Sin preámbulo ni \begin{document} — lo maneja main.tex
		%%%%%%%%%%%%%%%%%%%%%%%%%%%%%%%%%%%%%%%%%%%%%%%%%%%%%%%%%%%
		
		\section{Columna Sacra}
		
		\subsection{Proyección AP Oblicua (OPD / OPI) - Articulaciones Sacroilíacas}
		
		\subsubsection*{1. Identificación}
		
		\textbf{Tipo:} Proyección complementaria. Se realiza en ambas posiciones (OPD y OPI) para comparación bilateral.
		
		\textbf{Indicaciones clínicas principales:}
		\begin{itemize}
			\item Evaluación de luxación o subluxación de la articulación sacroilíaca.
			\item Procesos patológicos de la articulación sacroilíaca.
			\item Estudio comparativo bilateral de ambas articulaciones.
		\end{itemize}
		
		\subsubsection*{2. Técnica Base}
		
		\textbf{Receptor de imagen:} 24x30 cm. Orientación \textbf{longitudinal} (vertical). Se realizan ambas oblicuas.
		
		\textbf{DFRI:} 100 cm.
		
		\textbf{Técnica:} Básica para pacientes de gran espesor; intermedia para pacientes estándar; homogeneizada para pacientes muy delgados o pediátricos.
		
		\textbf{Bucky:} Bucky de mesa con grilla.
		
		\subsubsection*{3. Posicionamiento}
		
		\paragraph{Paciente}
		Decúbito supino. La cabeza se eleva sobre una almohada para mayor comodidad. Se eleva el lado en estudio aproximadamente 25--30°, apoyando el hombro, la parte inferior del tórax y la parte superior del muslo para mantener la oblicuidad.
		
		
		\begin{figure}[h]
			\centering
			\includegraphics[width=0.7\linewidth]{imagenes/col50}
			\caption{AP OPD}
			\label{fig:col50}
		\end{figure}
		
		
		\paragraph{Región de interés}
		Se alinea el cuerpo de forma que un plano sagital que pase 2,5 cm medial a la EIAS del lado elevado quede centrado sobre la línea media de la rejilla. Se verifica la rotación en varios puntos a lo largo de la espalda. El lado en estudio es el más alejado del receptor de imagen: posición OPI para demostrar la articulación derecha; posición OPD para demostrar la articulación izquierda.
		
		\paragraph{Rayo Central}
		Perpendicular al plano del receptor. Punto de entrada: \textbf{2,5 cm medial a la EIAS del lado elevado}.
		
		\begin{figure}[h]
			\centering
			\includegraphics[width=0.7\linewidth]{imagenes/col51}
			\caption{}
			\label{fig:col51}
		\end{figure}
		
		
		\paragraph{Receptor de imagen}
		Centrado a la altura de la EIAS. Debe incluir la articulación sacroilíaca del lado en estudio con márgenes adecuados para su evaluación.
		
		\paragraph{Respiración}
		Suspendida durante la exposición.
		
		\subsubsection*{4. Anatomía Visible}
		
		Se visualiza la articulación sacroilíaca del lado más alejado del receptor (lado elevado). La oblicuidad proyecta la articulación de forma que solo se radiografía la porción más posterior de la misma; el rayo central no pasa tangencial a la parte anterior de la articulación. Esto es consecuencia de la oblicuidad propia de la articulación sacroilíaca, cuya orientación no es completamente coronal, sino que presenta un componente posterior y otro anterior.
		
		\subsubsection*{5. Criterios de Evaluación}
		
		\textbf{Posicionamiento correcto:} Espacio articular sacroilíaco abierto con mínima superposición del ilíaco y el sacro. La articulación debe quedar centrada en la radiografía.
		
		\textbf{Ausencia de rotación:} Oblicuidad de 25--30° correctamente mantenida, verificada por la apertura del espacio articular sacroilíaco sin exceso de superposición ósea.
		
		\textbf{Penetración y contraste:} Técnica adecuada al biotipo del paciente (básica, intermedia u homogeneizada), con visualización nítida de la interlínea articular sacroilíaca.
		
		\subsubsection*{6. Puntos Críticos}
		
		\textbf{Confusión de lado:} Recordar que el lado en estudio es el elevado (más alejado del RI). OPI demuestra la articulación derecha; OPD demuestra la izquierda.
		
		\textbf{Visualización parcial de la articulación:} Por la oblicuidad anatómica de la articulación sacroilíaca, solo se radiografía su porción posterior. No es posible obtener una imagen completa en una sola proyección oblicua, lo que limita la valoración. Este enfoque no es de uso común por esta razón.
		
		\textbf{Adaptaciones:} En pacientes con limitación para el decúbito supino, ajustar el apoyo con cuñas o almohadillas para mantener los 25--30° de oblicuidad.
		
		\textbf{Trucos:} Siempre explorar ambos lados para comparación bilateral. La radioprotección gonadal es de aplicación obligatoria en varones mediante colimación estrecha centrada en la articulación; en mujeres es técnicamente difícil de aplicar sin comprometer la región de interés.
		
		
		\onecolumn
		\clearpage
		
		\begin{figure}
			\centering
			\includegraphics[width=0.7\linewidth]{imagenes/col52}
			\caption{}
			\label{fig:col52}
		\end{figure}
		
		
		\twocolumn
		\clearpage
		%----------------------------------------------------------
		%  Para agregar otro enfoque duplicá el bloque
		%  desde \subsection hasta \textbf{Trucos}
		%----------------------------------------------------------
		
		
	%%%%%%%%%%%%%%%%%%%%%%%%%%%%%%%%%%%%%%%%%%%%%%%%%%%%%%%%%%%
	%  SECCIÓN: Articulaciones Sacroilíacas
	%  Archivo: secciones/sacroiliacas_oblicua.tex
	%  Sin preámbulo ni \begin{document} — lo maneja main.tex
		%%%%%%%%%%%%%%%%%%%%%%%%%%%%%%%%%%%%%%%%%%%%%%%%%%%%%%%%%%%
		
		\section{Columna Sacra}
		
		\subsection{Proyección AP Oblicua (OPD / OPI) - Articulaciones Sacroilíacas}
		
		\subsubsection*{1. Identificación}
		
		\textbf{Tipo:} Proyección complementaria. Se realiza en ambas posiciones (OPD y OPI) para comparación bilateral.
		
		\textbf{Indicaciones clínicas principales:}
		\begin{itemize}
			\item Evaluación de luxación o subluxación de la articulación sacroilíaca.
			\item Procesos patológicos de la articulación sacroilíaca.
			\item Estudio comparativo bilateral de ambas articulaciones.
		\end{itemize}
		
		\subsubsection*{2. Técnica Base}
		
		\textbf{Receptor de imagen:} 24x30 cm. Orientación \textbf{longitudinal} (vertical). Se realizan ambas oblicuas.
		
		\textbf{DFRI:} 100 cm.
		
		\textbf{Técnica:} Básica para pacientes de gran espesor; intermedia para pacientes estándar; homogeneizada para pacientes muy delgados o pediátricos.
		
		\textbf{Bucky:} Bucky de mesa con grilla.
		
		\subsubsection*{3. Posicionamiento}
		
		\paragraph{Paciente}
		Decúbito supino. La cabeza se eleva sobre una almohada para mayor comodidad. Se eleva el lado en estudio aproximadamente 25--30°, apoyando el hombro, la parte inferior del tórax y la parte superior del muslo para mantener la oblicuidad.
		
		\begin{figure}[h]
			\centering
			\includegraphics[width=0.7\linewidth]{imagenes/col53}
			\caption{}
			\label{fig:col53}
		\end{figure}
		
		
		\paragraph{Región de interés}
		Se alinea el cuerpo de forma que un plano sagital que pase 2,5 cm medial a la EIAS del lado elevado quede centrado sobre la línea media de la rejilla. Se verifica la rotación en varios puntos a lo largo de la espalda. El lado en estudio es el más alejado del receptor de imagen: posición OPI para demostrar la articulación derecha; posición OPD para demostrar la articulación izquierda.
		
		\paragraph{Rayo Central}
		Perpendicular al plano del receptor. Punto de entrada: \textbf{2,5 cm medial a la EIAS del lado elevado}.
		
		\begin{figure}[h]
			\centering
			\includegraphics[width=0.7\linewidth]{imagenes/col54}
			\caption{}
			\label{fig:col54}
		\end{figure}
		
		
		\paragraph{Receptor de imagen}
		Centrado a la altura de la EIAS. Debe incluir la articulación sacroilíaca del lado en estudio con márgenes adecuados para su evaluación.
		
		\paragraph{Respiración}
		Suspendida durante la exposición.
		
		\subsubsection*{4. Anatomía Visible}
		
		Se visualiza la articulación sacroilíaca del lado más alejado del receptor (lado elevado). La oblicuidad proyecta la articulación de forma que solo se radiografía la porción más posterior de la misma; el rayo central no pasa tangencial a la parte anterior de la articulación. Esto es consecuencia de la oblicuidad propia de la articulación sacroilíaca, cuya orientación no es completamente coronal, sino que presenta un componente posterior y otro anterior.
		
		\subsubsection*{5. Criterios de Evaluación}
		
		\textbf{Posicionamiento correcto:} Espacio articular sacroilíaco abierto con mínima superposición del ilíaco y el sacro. La articulación debe quedar centrada en la radiografía.
		
		\textbf{Ausencia de rotación:} Oblicuidad de 25--30° correctamente mantenida, verificada por la apertura del espacio articular sacroilíaco sin exceso de superposición ósea.
		
		\textbf{Penetración y contraste:} Técnica adecuada al biotipo del paciente (básica, intermedia u homogeneizada), con visualización nítida de la interlínea articular sacroilíaca.
		
		\subsubsection*{6. Puntos Críticos}
		
		\textbf{Confusión de lado:} Recordar que el lado en estudio es el elevado (más alejado del RI). OPI demuestra la articulación derecha; OPD demuestra la izquierda.
		
		\textbf{Visualización parcial de la articulación:} Por la oblicuidad anatómica de la articulación sacroilíaca, solo se radiografía su porción posterior. No es posible obtener una imagen completa en una sola proyección oblicua, lo que limita la valoración. Este enfoque no es de uso común por esta razón.
		
		\textbf{Adaptaciones:} En pacientes con limitación para el decúbito supino, ajustar el apoyo con cuñas o almohadillas para mantener los 25--30° de oblicuidad.
		
		\textbf{Trucos:} Siempre explorar ambos lados para comparación bilateral. La radioprotección gonadal es de aplicación obligatoria en varones mediante colimación estrecha centrada en la articulación; en mujeres es técnicamente difícil de aplicar sin comprometer la región de interés.
		
		\subsection{Proyección PA Oblicua (OAD / OAI) --- Articulaciones Sacroilíacas}
		
		\subsubsection*{1. Identificación}
		
		\textbf{Tipo:} Proyección complementaria. Se realiza en ambas posiciones (OAD y OAI) para comparación bilateral.
		
		\textbf{Indicaciones clínicas principales:}
		\begin{itemize}
			\item Traumatismo de la articulación sacroilíaca.
			\item Pinzamiento, luxación o subluxación de la articulación sacroilíaca.
			\item Procesos patológicos de la articulación sacroilíaca.
		\end{itemize}
		
		\subsubsection*{2. Técnica Base}
		
		\textbf{Receptor de imagen:} 18x24 cm o 24x30 cm. Orientación \textbf{longitudinal} (vertical).
		
		\textbf{DFRI:} 100 cm.
		
		\textbf{Técnica:} Básica para pacientes de gran espesor; intermedia para pacientes estándar; homogeneizada para pacientes pediátricos o muy delgados.
		
		\textbf{Bucky:} Bucky de mesa con grilla.
		
		\subsubsection*{3. Posicionamiento}
		
		\paragraph{Paciente}
		Posición semiprona. El lado a estudiar queda más próximo al receptor de imagen. Se coloca almohada bajo la cabeza. El paciente se apoya sobre su antebrazo y flexiona la rodilla del lado elevado para estabilizar la posición.
		
		\paragraph{Región de interés}
		Se rota al paciente del lado de interés hacia la mesa unos 25--50°. Se ajusta el cuerpo de forma que su eje longitudinal sea paralelo al eje longitudinal de la mesa. Se centra el cuerpo de forma que un punto 2,5 cm medial a la EIAS más próxima al RI quede centrado sobre la línea media de la rejilla. OAD para demostrar la articulación derecha; OAI para demostrar la articulación izquierda.
		
		\paragraph{Rayo Central}
		Perpendicular al plano del receptor. Punto de entrada: \textbf{2,5 cm medial a la EIAS más próxima al RI}.

		
		\paragraph{Receptor de imagen}
		Centrado a la altura de la EIAS. Debe incluir la articulación sacroilíaca del lado en estudio.
		
		\paragraph{Respiración}
		Suspendida durante la exposición.
		
		\subsubsection*{4. Anatomía Visible}
		
		Se visualiza la articulación sacroilíaca del lado más próximo al receptor (lado en estudio). La colimación debe ser estrecha, limitada a la articulación sacroilíaca de interés.
		
		\subsubsection*{5. Criterios de Evaluación}
		
		\textbf{Posicionamiento correcto:} Articulación sacroilíaca del lado en estudio centrada en la radiografía, con apertura adecuada del espacio articular.
		
		\textbf{Ausencia de rotación:} Oblicuidad de 25--50° correctamente mantenida, verificada por la visualización del espacio articular sin superposición excesiva.
		
		\textbf{Penetración y contraste:} Técnica adecuada al biotipo del paciente (básica, intermedia u homogeneizada), con nitidez de la interlínea articular sacroilíaca.
		
		\subsubsection*{6. Puntos Críticos}
		
		\textbf{Confusión de lado:} En esta proyección el lado en estudio es el más próximo al RI (a diferencia de la AP oblicua). OAD demuestra la articulación derecha; OAI demuestra la izquierda.
		
		\textbf{Ángulo insuficiente o excesivo:} El rango de oblicuidad es amplio (25--50°); un ángulo inadecuado genera superposición del ilíaco sobre el sacro y cierre del espacio articular.
		
		\textbf{Adaptaciones:} En pacientes con dificultad para mantener la posición semiprona, utilizar cuñas o almohadillas bajo el tórax y la pelvis para sostener la oblicuidad requerida.
		
		\textbf{Trucos:} Siempre realizar ambas posiciones (OAD y OAI) para comparación bilateral. La colimación estrecha es obligatoria en todos los pacientes. Confirmar la oblicuidad palpando la posición de la EIAS antes de exponer.
		
		
		
		\onecolumn
		\clearpage
		
		\begin{figure}
			\centering
			\includegraphics[width=0.7\linewidth]{imagenes/col55}
			\caption{OAI con 20° de angulación caudal}
			\label{fig:col55}
		\end{figure}
		
		
		
		\twocolumn
		\clearpage
		
		%----------------------------------------------------------
		%  Para agregar otro enfoque duplicá el bloque
		%  desde \subsection hasta \textbf{Trucos}
		%----------------------------------------------------------
		
		
%%%%%%%%%%%%%%%%%%%%%%%%%%%%%%%%%%%%%%%%%%%%%%%%%%%%%%%%%%%
%  SECCIÓN: Sacro y Cóccix
%  Archivo: secciones/sacro_coccix.tex
%  Sin preámbulo ni \begin{document} — lo maneja main.tex
	%%%%%%%%%%%%%%%%%%%%%%%%%%%%%%%%%%%%%%%%%%%%%%%%%%%%%%%%%%%
	
	\section{Columna Sacra}
	
	% ----------------------------------------------------------
	%  ENFOQUE 1: AP / PA AXIAL — SACRO
	% ----------------------------------------------------------
	
	\subsection{Proyección AP / PA Axial - Sacro}
	
	\subsubsection*{1. Identificación}
	
	\textbf{Tipo:} Proyección básica. La variante PA se emplea como alternativa posicional cuando la AP no es viable.
	
	\textbf{Indicaciones clínicas principales:}
	\begin{itemize}
		\item Patologías del sacro
		\item Fracturas o lesiones destructivas del sacro
		\item Evaluación de agujeros sacros, alas sacras y articulación lumbosacra
	\end{itemize}
	
	\subsubsection*{2. Técnica Base}
	
	\textbf{Receptor de imagen:} 24x30 cm.
	
	\textbf{DFRI:} 100 cm.
	
	\textbf{Técnica:} Intermedia. Permite equilibrar la visualización de las estructuras óseas del sacro con adecuado contraste.
	
	\textbf{Bucky:} Bucky de mesa con grilla.
	
	\subsubsection*{3. Posicionamiento}
	
	\paragraph{Paciente}
	\textbf{AP:} Decúbito supino. \textbf{PA:} Decúbito prono, empleada en pacientes con lesión dolorosa o enfermedad destructiva, sin pérdida de detalle diagnóstico respecto a la AP.
	En ambos casos el plano medio sagital perpendicular y centrado al receptor.
	
	% TODO: \usepackage{graphicx} required
	\begin{figure}[h]
		\centering
		\includegraphics[width=0.7\linewidth]{imagenes/col56}
		\caption{}
		\label{fig:col56}
	\end{figure}
	
	
	\paragraph{Región de interés}
	Ambas EIAS equidistantes de la grilla para evitar rotación. Codos flexionados, brazos en posición simétrica bilateral cómoda. En supino, apoyo bajo las rodillas para reducir la lordosis lumbar y acercar el sacro al receptor. Los huesos deben estar tan próximos al receptor como sea posible.
	
	\paragraph{Rayo Central} - punto de entrada 5 cm por encima de la sínfisis púbica.\\
	\textbf{PA (prono):} 15° caudal - punto de entrada centrado sobre la curvatura sacra.\\
	La dirección se invierte al cambiar de posición para lograr el mismo efecto de desplegado del sacro.
	
	\paragraph{Receptor de imagen}
	Centrado de modo que incluya desde L5 hasta el cóccix.
	
	\paragraph{Respiración}
	Suspendida durante la exposición.
	
	\subsubsection*{4. Anatomía Visible}
	
	Sacro con su curvatura rectificada gracias a la angulación del RC. Se observan los agujeros sacros anteriores, las alas sacras simétricas, la articulación lumbosacra (L5--S1) y los huesos pubianos libres de superposición sobre el sacro. La angulación de 15° compensa la inclinación anterior del sacro, desplegando su curvatura.
	
	\subsubsection*{5. Criterios de Evaluación}
	
	\textbf{Posicionamiento correcto:} Sacro centrado, libre de acortamiento y con curvatura rectificada. Huesos pubianos no superpuestos al sacro. Ausencia de materia fecal superpuesta a la región de interés.
	
	\textbf{Ausencia de rotación:} Alas sacras simétricas en ambos lados. Ambas EIAS equidistantes de la línea media.
	
	\textbf{Penetración y contraste:} Escala de contraste corta que permita diferenciar las estructuras óseas del sacro.
	
	\subsubsection*{6. Puntos Críticos}
	
	\textbf{Angulación insuficiente del RC:} Provoca superposición del pubis sobre el sacro. El ángulo exacto (entre 10° y 20°) se ajusta según el promontorio y la morfología de caderas del paciente.
	
	\textbf{Confusión en la dirección de angulación al pasar a prono:} En PA la angulación es caudal, opuesta a la AP. Invertir la dirección es el error más frecuente al cambiar de posición.
	
	\textbf{Adaptaciones:} En pacientes con lesión dolorosa o enfermedad destructiva, emplear la variante PA (prono) con angulación caudal de 15°.
	
	\textbf{Trucos:} El sacro visto de perfil tiende a irse hacia adelante. La angulación del RC compensa esa inclinación natural; el ángulo exacto varía por paciente según el promontorio y la morfología de las caderas.
	
	
	\clearpage
	\onecolumn
	
	% TODO: \usepackage{graphicx} required
	\begin{figure}
		\centering
		\includegraphics[width=0.7\linewidth]{imagenes/col57}
		\caption{AP de Sacro}
		\label{fig:col57}
	\end{figure}
			
	\twocolumn
	\clearpage
	
		\clearpage
	\onecolumn
	
	% TODO: \usepackage{graphicx} required
	\begin{figure}
		\centering
		\includegraphics[width=0.7\linewidth]{imagenes/col58}
		\caption{AP de Sacro}
		\label{fig:col58}
	\end{figure}
	
	\twocolumn
	\clearpage
	% ----------------------------------------------------------
	%  ENFOQUE 2: AP / PA AXIAL — CÓCCIX
	% ----------------------------------------------------------
	
	\subsection{Proyección AP / PA Axial - Cóccix}
	
	\subsubsection*{1. Identificación}
	
	\textbf{Tipo:} Proyección básica. La variante PA se emplea como alternativa posicional cuando la AP no es viable.
	
	\textbf{Indicaciones clínicas principales:}
	\begin{itemize}
		\item Patologías del cóccix
		\item Traumatismos y fracturas coccígeas
		\item Evaluación de la alineación de los segmentos del cóccix
	\end{itemize}
	
	\subsubsection*{2. Técnica Base}
	
	\textbf{Receptor de imagen:} 18x24 cm.
	
	\textbf{DFRI:} 100 cm.
	
	\textbf{Técnica:} Intermedia.
	
	\textbf{Bucky:} Bucky de mesa con grilla.
	
	\subsubsection*{3. Posicionamiento}
	
	\paragraph{Paciente}
	\textbf{AP:} Decúbito supino. \textbf{PA:} Decúbito prono, empleada en pacientes con lesión dolorosa o enfermedad destructiva, sin pérdida de detalle diagnóstico respecto a la AP.
	En ambos casos el plano medio sagital perpendicular y centrado al receptor.
	
	% TODO: \usepackage{graphicx} required
	\begin{figure}[h]
		\centering
		\includegraphics[width=0.7\linewidth]{imagenes/col59}
		\caption{}
		\label{fig:col59}
	\end{figure}
	
	
	\paragraph{Región de interés}
	Ambas EIAS equidistantes de la grilla. Codos flexionados, brazos en posición simétrica bilateral cómoda. En supino, apoyo bajo las rodillas. Los huesos deben estar tan próximos al receptor como sea posible.
	
	\paragraph{Rayo Central}
	\textbf{AP (supino):} 10° caudal - punto de entrada 5 cm por encima de la sínfisis púbica.\\
	\textbf{PA (prono):} 10° craneal - punto de entrada centrado sobre el cóccix prominente.\\
	La dirección se invierte al cambiar de posición para lograr el mismo efecto de desplegado de los segmentos coccígeos.
	
	\paragraph{Receptor de imagen}
	Centrado sobre la región coccígea. Colimación estrecha para mejorar la visualización y reducir la radiación dispersa. Incluir desde el sacro distal hasta el extremo del cóccix.
	
	\paragraph{Respiración}
	Suspendida durante la exposición.
	
	\subsubsection*{4. Anatomía Visible}
	
	Segmentos del cóccix sin superposición entre sí, centrados en la imagen, y articulación sacrococcígea. El cóccix presenta una curvatura cifótica opuesta a la del sacro, por lo que la angulación del RC en AP es caudal (inversa a la del sacro). La colimación estrecha mejora el contraste al reducir la radiación dispersa.
	
	\subsubsection*{5. Criterios de Evaluación}
	
	\textbf{Posicionamiento correcto:} Cóccix centrado y visible en su totalidad. Segmentos del cóccix no superpuestos. Colimación estrecha evidente en la imagen.
	
	\textbf{Ausencia de rotación:} Cóccix alineado sobre la línea media sin desviación lateral de sus segmentos.
	
	\textbf{Penetración y contraste:} Escala de contraste corta. La colimación estrecha mejora el contraste de la región.
	
	\subsubsection*{6. Puntos Críticos}
	
	\textbf{Confusión en la dirección de angulación respecto al sacro:} A diferencia del sacro (AP craneal), en el cóccix la AP es caudal por su curvatura cifótica opuesta. Es el error conceptual más frecuente al estudiar ambas regiones juntas.
	
	\textbf{Colimación inadecuada:} Una colimación amplia incrementa la radiación dispersa y degrada el contraste. Colimar estrictamente a la región de interés.
	
	\textbf{Adaptaciones:} En pacientes con lesión dolorosa, emplear la variante PA (prono) con angulación craneal de 10°.
	
	\textbf{Trucos:} El cóccix de perfil presenta una curvatura en cifosis, angulación contraria a la del sacro. La protección gonadal no puede colocarse en mujeres en esta proyección; en hombres se protegen las gónadas colocando el protector en los hombros.
	
	\onecolumn
	\clearpage
	
	% TODO: \usepackage{graphicx} required
	\begin{figure}
		\centering
		\includegraphics[width=0.7\linewidth]{imagenes/col60}
		\caption{AP Coxis}
		\label{fig:col60}
	\end{figure}
	
	
	\twocolumn
	\clearpage
	
	%----------------------------------------------------------
	%  Para agregar otro enfoque duplicá el bloque
	%  desde \subsection hasta \textbf{Trucos}
	%----------------------------------------------------------
	
	
%%%%%%%%%%%%%%%%%%%%%%%%%%%%%%%%%%%%%%%%%%%%%%%%%%%%%%%%%%%
%  SECCIÓN: Columna Lumbosacra — Sacro y Cóccix
%  Archivo: secciones/lateral_sacro_coccix.tex
%  Sin preámbulo ni \begin{document} — lo maneja main.tex
	%%%%%%%%%%%%%%%%%%%%%%%%%%%%%%%%%%%%%%%%%%%%%%%%%%%%%%%%%%%
	
	\section{Columna Coxis}
	
	\subsection{Proyección Lateral: Sacro y Cóccix}
	
	\subsubsection*{1. Identificación}
	
	\textbf{Tipo:} Proyección complementaria. Forma par con la proyección AP de sacro y cóccix.
	
	\textbf{Indicaciones clínicas principales:}
	\begin{itemize}
		\item Patologías del sacro
		\item Patologías del cóccix
	\end{itemize}
	
	\subsubsection*{2. Técnica Base}
	
	\textbf{Receptor de imagen (sacro):} 24$\times$30 cm, centrado en la línea media del sacro.
	
	\textbf{Receptor de imagen (cóccix):} 18$\times$24 cm, centrado en la línea media del cóccix.
	
	\textbf{DFRI:} 100 cm.
	
	\textbf{Técnica:} Intermedia.
	
	\textbf{Bucky:} Sí.
	
	\subsubsection*{3. Posicionamiento}
	
	\paragraph{Paciente}
	Paciente en decúbito lateral, girado hacia el lado indicado. Flexionar caderas y rodillas en una posición cómoda. Ajustar brazos en una posición cómoda. Colocar soporte bajo el cuerpo para situar el eje longitudinal de la columna horizontal. Ajustar pelvis y hombros para obtener una posición lateral verdadera.
	
	% TODO: \usepackage{graphicx} required
	\begin{figure}[h]
		\centering
		\includegraphics[width=0.7\linewidth]{imagenes/col61}
		\caption{Lateral de Coxis}
		\label{fig:col61}
	\end{figure}
	
	% TODO: \usepackage{graphicx} required
	\begin{figure}[h]
		\centering
		\includegraphics[width=0.7\linewidth]{imagenes/col62}
		\caption{Lateral de Sacro}
		\label{fig:col62}
	\end{figure}
	
	
	\paragraph{Región de interés}
	La posición exacta del sacro depende de la curvatura pélvica del paciente. Los márgenes posteriores de los isquiones y los iliacos deben quedar casi superpuestos entre sí. Colimación estrecha sobre la región de interés.
	
	\paragraph{Rayo Central — Sacro}
	RC perpendicular al plano del receptor. Punto de entrada: \textbf{EIAS, dirigido 9 cm posterior a la misma}.
	
	\paragraph{Rayo Central — Cóccix}
	RC perpendicular al plano del receptor. Punto de entrada: \textbf{9 cm posterior y 5 cm inferior a la EIAS}.
	
	\paragraph{Respiración}
	Suspendida durante la exposición.
	
	\subsubsection*{4. Anatomía Visible}
	
	Sacro y cóccix en proyección lateral. La superposición de los márgenes posteriores de los isquiones e iliacos entre sí confirma el posicionamiento lateral verdadero. La curvatura natural del sacro queda demostrada sin rotación pélvica.
	
	\subsubsection*{5. Criterios de Evaluación}
	
	\textbf{Posicionamiento correcto:} Márgenes posteriores de los isquiones y los iliacos casi superpuestos entre sí; eje longitudinal de la columna horizontal al receptor.
	
	\textbf{Ausencia de rotación:} Superposición simétrica de los elementos posteriores de la pelvis; pelvis y hombros en lateral verdadero sin inclinación.
	
	\textbf{Penetración y contraste:} Técnica intermedia que permita visualizar el sacro y el cóccix con densidad adecuada, diferenciando las estructuras óseas del tejido blando circundante.
	
	\subsubsection*{6. Puntos Críticos}
	
	\textbf{Rotación pélvica:} Genera superposición asimétrica de los isquiones e iliacos, distorsionando la imagen del sacro y el cóccix; corregir ajustando pelvis y hombros a lateral verdadero.
	
	\textbf{Eje de columna no horizontal:} Si el soporte bajo el cuerpo es insuficiente, la columna queda inclinada y el RC no incide perpendicularmente sobre la región; verificar la horizontalidad antes de la exposición.
	
	\textbf{Adaptaciones:} En pacientes con limitaciones de movilidad o dolor agudo, ajustar el grado de flexión de caderas y rodillas según tolerancia, manteniendo la prioridad de lograr el lateral verdadero.
	
	\textbf{Trucos:} La EIAS es el repere principal y es fácilmente palpable con el paciente en decúbito lateral. Para el sacro, medir 9 cm posterior a la EIAS; para el cóccix, combinar 9 cm posterior y 5 cm inferior. Solicitar colimación estrecha para reducir radiación dispersa y mejorar el contraste en ambas estructuras.
	
	
	\onecolumn
	\clearpage
	
	% TODO: \usepackage{graphicx} required
	\begin{figure}
		\centering
		\includegraphics[width=1\linewidth]{imagenes/col63}
		\caption{\textbf{A}. Proyección lateral del sacro. \textbf{B}. Proyección lateral del cóccix.
		}
		\label{fig:col63}
	\end{figure}
		
	
	\twocolumn
	\clearpage
	
	%----------------------------------------------------------
	%  Para agregar otro enfoque duplicá el bloque
	%  desde \subsection hasta \textbf{Trucos}
	%----------------------------------------------------------
	
	
%%%%%%%%%%%%%%%%%%%%%%%%%%%%%%%%%%%%%%%%%%%%%%%%%%%%%%%%%%%
%  SECCIÓN: Diagnóstico por Imágenes — Escoliosis
%  Archivo: secciones/dc_escoliosis.tex
%  Sin preámbulo ni \begin{document} — lo maneja main.tex
	%%%%%%%%%%%%%%%%%%%%%%%%%%%%%%%%%%%%%%%%%%%%%%%%%%%%%%%%%%%
	
	\section{Columna - Escoliosis}
	
	%----------------------------------------------------------
	\subsection{Definición y Características Generales}
	
	\subsubsection*{Concepto}
	
	La escoliosis es la desviación de la columna vertebral con deformidad tridimensional. La deformidad predominante se localiza en el \textbf{plano coronal} (desviación derecha–izquierda), con un componente rotacional asociado que puede afectar la cifosis dorsal y la lordosis lumbar. Para ser considerada escoliosis, la angulación debe superar los \textbf{10°}.
	
	\subsubsection*{Planos de la Columna}
	
	\textbf{Plano sagital:} contiene las curvaturas fisiológicas normales.
	\begin{itemize}
		\item Lordosis cervical y lumbar
		\item Cifosis torácica y sacra
	\end{itemize}
	
	\textbf{Plano coronal:} toda curvatura en este plano es \textbf{patológica} (desviación lateral derecha–izquierda o izquierda–derecha).
	
	
	\begin{figure}[h]
		\centering
		\includegraphics[width=0.7\linewidth]{imagenes/col64}
		\caption{}
		\label{fig:col64}
	\end{figure}
	
	
	%----------------------------------------------------------
	\subsection{Vértebras de Referencia en Escoliosis}
	
		% TODO: \usepackage{graphicx} required
	\begin{figure}[h]
		\centering
		\includegraphics[width=0.7\linewidth]{imagenes/col67}
		\caption{}
		\label{fig:col65}
	\end{figure}
	
	
	\subsubsection*{Vértebra Proximal (Extremo Superior)}
	Es la vértebra más superior de la curvatura. Su superficie superior gira en sentido máximo hacia la concavidad de la curva. Es la primera vértebra en manifestarse en la curvatura.
	
	\subsubsection*{Vértebra Distal (Terminal)}
	Es la vértebra más inferior de la curvatura. Su superficie inferior gira en sentido máximo hacia la concavidad de la curva. Es la última vértebra en manifestarse en la curvatura.
	
	\subsubsection*{Vértebra Ápex (Apical)}
	Es la vértebra más desviada del eje vertical del paciente. Es también la más rotada, lo que determina su mayor desplazamiento lateral. Corresponde al punto máximo de la curvatura.
	
	
	%----------------------------------------------------------
	\subsection{Clasificación de las Curvas}
	

	
	\subsubsection*{Curva Primaria}
	La curva principal de la escoliosis. Presenta mayor rotación y desviación de los cuerpos vertebrales. Es la menos flexible: al colocar un cajón bajo el pie del lado convexo, esta curva no se modifica. Corresponde siempre a la curva con \textbf{mayor ángulo de Cobb}.
	
	\subsubsection*{Curva Secundaria (Compensadora)}
	Son las curvas que el cuerpo genera para compensar la curva primaria, en dirección opuesta. Son menos acentuadas y más flexibles. Siempre están presentes junto a la curva primaria. El ángulo de la curva compensadora es \textbf{menor} que el de la curva primaria.
	
	\textbf{Ejemplo:} curva primaria dorsal + curva compensadora lumbar.
	
	\subsubsection*{Curva Estructural}
	La curva está fija en ese segmento espinal. No puede corregirse por completo en radiografías con el paciente en decúbito dorsal ni en flexión lateral. Se perdió flexibilidad en la columna.
	
	\subsubsection*{Curva No Estructural}
	No tiene rotación fija ni angulación lateral permanente. Puede corregirse completamente por tracción o flexión lateral. Corresponde a escoliosis postural o por mala posición; al colocar un cajoncito, esta curva mejora.
	
	\onecolumn
	
	\begin{figure}
		\centering
		\includegraphics[width=0.6\linewidth]{imagenes/col65}
		\caption{}
		\label{fig:col65}
	\end{figure}
	
	\twocolumn	
	%----------------------------------------------------------
	\clearpage
	\subsection{Ángulo de Cobb}
	
	\subsubsection*{Técnica de Medición}
	
	\paragraph{Paso 1}
	Trazar una línea paralela al \textbf{platillo superior} de la vértebra límite superior de la curvatura.
	
	\paragraph{Paso 2}
	Trazar una línea paralela al \textbf{platillo inferior} de la vértebra límite inferior de la curvatura.
	
	\paragraph{Paso 3}
	Trazar perpendiculares a ambas líneas y medir el ángulo formado entre ellas.
	
	\paragraph{Referencia alternativa}
	Si los platillos no son identificables en la radiografía, se toman como referencia los \textbf{bordes de los pedículos}.
	
	\subsubsection*{Interpretación Clínica}
	
	\textbf{Menos de 20°:} observación y control periódico.
	
	\textbf{Entre 20° y 40°:} tratamiento ortopédico.
	
	\textbf{Mayor de 50°:} indicación quirúrgica.
	
	\subsubsection*{Utilidad del Ángulo de Cobb}
	Cuantifica la desviación vertebral en el plano coronal. Permite diferenciar la \textbf{curva mayor} (mayor ángulo) de las \textbf{curvas menores} (ángulo menor). Las curvaturas menores tienden a compensar la curvatura mayor. Siempre se visualizan al menos dos curvaturas: una escoliótica y una compensatoria.
	
	\clearpage
	\onecolumn
	
	% TODO: \usepackage{graphicx} required
	\begin{figure}
		\centering
		\includegraphics[width=0.7\linewidth]{imagenes/col66}
		\caption{}
		\label{fig:col66}
	\end{figure}
	
	
	\twocolumn
	\clearpage
	
	%----------------------------------------------------------
\subsection{Signo de Risser}

\subsubsection*{Definición}
Indicador radiográfico de \textbf{madurez esquelética} desarrollado por el Dr. Joseph C. Risser en 1958. Evalúa el grado de osificación de la \textbf{apófisis ilíaca (cresta ilíaca)} en radiografías de pelvis. La osificación de las apófisis ilíacas ocurre de manera paralela a la de las apófisis vertebrales, permitiendo estimar el potencial de crecimiento espinal restante.

\subsubsection*{Proceso de Osificación}
La osificación comienza en la región \textbf{anterolateral} de la cresta ilíaca (cerca de la espina ilíaca anterosuperior), avanza medialmente hacia la espina ilíaca posterosuperior y finalmente se fusiona en dirección medial a lateral.

\subsubsection*{Escala de Clasificación (0--5)}

\textbf{Risser 0:} Sin osificación visible. Crecimiento significativo restante. Alto riesgo de progresión de escoliosis. Etapa prepuberal o pubertad temprana.

\textbf{Risser 1:} Osificación del 25\% de la cresta ilíaca. Potencial de crecimiento considerable. Etapa prepuberal o pubertad temprana.

\textbf{Risser 2:} Osificación del 50\%. Potencial de crecimiento moderado. Antes o durante el pico de crecimiento.

\textbf{Risser 3:} Osificación del 75\%. Potencial limitado. Etapa de desaceleración del crecimiento.

\textbf{Risser 4:} Osificación completa (100\%) sin fusión. Potencial mínimo. Casi cesación del crecimiento.

\textbf{Risser 5:} Fusión completa de la apófisis con el ilion. Madurez esquelética completa. Sin crecimiento restante.

\subsubsection*{Importancia Clínica}

\textbf{Predicción de progresión:} Pacientes con Risser 0--2 presentan mayor riesgo de progresión de la curvatura escoliótica y requieren seguimiento más intensivo.

\textbf{Decisión terapéutica:} El corsé ortopédico es más efectivo en grados bajos (0--2). A mayor madurez, mayor probabilidad de indicación quirúrgica. Pacientes con Risser 4--5 requieren menos seguimiento activo.

\subsubsection*{Importancia Técnica --- Colimación}
Al dejar el colimador abierto en la radiografía panorámica de columna, se incluyen la pelvis y las crestas ilíacas en una sola toma. Esto \textbf{evita una radiografía adicional de pelvis} y reduce la dosis de radiación al paciente.

\subsubsection*{Técnica Radiológica}
Se evalúa en la \textbf{radiografía AP de pelvis}, incluida en el estudio de columna completa, con el paciente en bipedestación. Se requiere visualización completa de \textbf{ambas crestas ilíacas} y evaluación bilateral del porcentaje de osificación y de la fusión apofisaria.
	
	\onecolumn
	
	
	% TODO: \usepackage{graphicx} required
	\begin{figure}
		\centering
		\includegraphics[width=0.8\linewidth]{imagenes/col71}
		\caption{}
		\label{fig:col70}
	\end{figure}
	
	\begin{figure}
		\centering
		\includegraphics[width=0.8\linewidth]{imagenes/col70}
		\caption{}
		\label{fig:col70}
	\end{figure}
	
	
	
	
	\twocolumn
	\clearpage
	%----------------------------------------------------------
	\subsection{Estudio Radiológico de la Escoliosis}
	
	\subsubsection*{Objetivos del Estudio}
	\begin{itemize}
		\item Visualizar las deformidades de los ejes de la columna vertebral.
		\item Cuantificar el grado de desviación mediante el ángulo de Cobb.
		\item Estimar el grado de madurez ósea del paciente (Signo de Risser).
	\end{itemize}
	
	\subsubsection*{Proyecciones Indicadas}
	\begin{itemize}
		\item Proyección panorámica de columna AP (PA) — frente.
		\item Proyección panorámica lateral de columna.
	\end{itemize}
	
	\subsubsection*{Elección del Perfil}
	La elección del perfil depende de la curvatura del paciente. Se coloca siempre la \textbf{concavidad de la curvatura hacia el lado del tubo}.
	
	\textbf{Ejemplo:} si la curvatura tiene concavidad izquierda → perfil derecho.
	
	\subsubsection*{Posición del Paciente — Aspectos Críticos}
	
	\textbf{Bipedestación:} el paciente debe estar de pie de forma \textbf{obligatoria}. La radiografía en carga refleja la posición anatómica real.
	
	\textbf{Descalzo:} el paciente debe estar sin calzado de forma \textbf{fundamental}. El calzado (especialmente con tacones o desgaste desigual) puede generar un falso diagnóstico al alterar la postura basal.
	
	\textbf{Distribución del peso:} el peso corporal debe estar distribuido equitativamente en ambos pies, para valorar la columna con la carga activa real del paciente.
	
	
	
	%----------------------------------------------------------
	%  Fin de sección dc_escoliosis.tex
	%----------------------------------------------------------
	\clearpage
	
%%%%%%%%%%%%%%%%%%%%%%%%%%%%%%%%%%%%%%%%%%%%%%%%%%%%%%%%%%%
%  SECCIÓN: Columna Vertebral — Proyecciones Panorámicas
%  Archivo: secciones/columna_panoramica.tex
%  Sin preámbulo ni \begin{document} — lo maneja main.tex
	%%%%%%%%%%%%%%%%%%%%%%%%%%%%%%%%%%%%%%%%%%%%%%%%%%%%%%%%%%%
	
	\section{Columna Escoliosis}
	
	\subsection{Panorámica de Columna AP/PA - Frente}
	
	\subsubsection*{1. Identificación}
	
	\textbf{Tipo:} Proyección básica. Forma par con la proyección panorámica de columna en perfil.
	
	\textbf{Indicaciones clínicas principales:}
	\begin{itemize}
		\item Escoliosis: evaluación de grado e intensidad de la curvatura.
		\item Caracterización de deformidades de los ejes de la columna vertebral.
		\item Estimación del grado de madurez esquelética (índice de Risser) como factor pronóstico para el manejo médico o quirúrgico.
	\end{itemize}
	
	\subsubsection*{2. Técnica Base}
	
	\textbf{Receptor de imagen:} 35x43 cm. Orientación \textbf{longitudinal} (vertical). Borde inferior ubicado 3--5 cm por debajo de la cresta ilíaca.
	
	\textbf{DFRI:} 180--200 cm (mínimo 150 cm). Se requiere mayor SID cuando se utiliza un receptor de mayor tamaño, a fin de obtener la colimación requerida.
	
	\textbf{Técnica:} Homogenizada en paciente pediátrico (80\,kVp --- 40\,mAs; perfil: 80\,kVp --- 80\,mAs o 100\,kVp --- 100\,mAs). Técnica intermedia en adulto.
	
	\textbf{Bucky:} Bucky vertical con grilla.
	
	\subsubsection*{3. Posicionamiento}
	
	\paragraph{Paciente}
	De pie (bipedestación), descalzo, con el peso distribuido equitativamente sobre ambos miembros inferiores. Se recomienda la posición PA (pecho pegado al Bucky vertical) por razones de radioprotección. En pacientes pediátricos se prefiere AP, dado que los niños tienden a moverse y es más difícil mantenerlos en PA.
	
	% TODO: \usepackage{graphicx} required
	\begin{figure}[h]
		\centering
		\includegraphics[width=0.7\linewidth]{imagenes/col76}
		\caption{}
		\label{fig:col76}
	\end{figure}
	
	
	\paragraph{Región de interés}
	Brazos a los costados del cuerpo. Verificar ausencia de rotación del torso y de la pelvis. Se toma referencia desde el plano medio sagital, 4 dedos lateralmente. Nota: la escoliosis puede generar rotación vertebral inherente, por lo que cierta rotación es inevitable y esperada.
	
	\paragraph{Rayo Central}
	Perpendicular al plano del receptor. Punto de entrada: \textbf{punto medio del receptor de imagen}.
	
	\paragraph{Receptor de imagen}
	Centrado de modo que el borde inferior quede 3-5 cm por debajo de la cresta ilíaca. Debe incluir vértebras torácicas, lumbares y quirúrgicas, y la totalidad de las crestas ilíacas (especialmente en adolescentes). Colimación a los cuatro lados; no se recomienda colimación estrecha, ya que es necesario evaluar las deformidades de costillas y pelvis adyacentes.
	
	\paragraph{Respiración}
	Suspendida al final de una espiración: se indica al paciente que inhale y luego aguante la respiración.
	
	\subsubsection*{4. Anatomía Visible}
	
	Vértebras torácicas, lumbares y quirúrgicas en proyección frontal. Totalidad de las crestas ilíacas, con al menos 2.5 cm visibles por encima de su borde superior en pacientes adolescentes, como indicador del índice de Risser. Regiones costales adyacentes y pelvis en sus márgenes. Puede evidenciarse rotación vertebral asociada a la curvatura escoliótica, lo cual es un hallazgo esperado y no un error técnico.
	
	\subsubsection*{5. Criterios de Evaluación}
	
	\textbf{Posicionamiento correcto:} Visualización de al menos 2.5 cm de crestas ilíacas por encima de su borde superior en pacientes adolescentes. Inclusión completa de vértebras torácicas, lumbares y quirúrgicas en un solo receptor.
	
	\textbf{Ausencia de rotación:} Simetría del torso y pelvis en la medida en que la patología lo permita. Se acepta rotación inherente a la escoliosis.
	
	\textbf{Penetración y contraste:} Adecuada visualización de los cuerpos vertebrales y espacios intervertebrales en toda la extensión de la columna. La técnica homogeneizada (niño) o intermedia (adulto) debe garantizar uniformidad de densidad en regiones torácica y lumbar simultáneamente.
	
	\subsubsection*{6. Puntos Críticos}
	
	\textbf{Colimación excesivamente estrecha:} Impide la evaluación de costillas adyacentes y pelvis, perdiendo información clínica relevante para la clasificación de la deformidad. Colimar a los cuatro lados sin estrechar lateralmente.
	
	\textbf{Omisión de crestas ilíacas:} En adolescentes, no visualizar al menos 2.5 cm de crestas ilíacas impide determinar el índice de Risser, dato esencial para la decisión médico-quirúrgica. Ajustar el borde inferior del receptor a 3--5 cm por debajo de la cresta.
	
	\textbf{Adaptaciones:} Si las estructuras no entran en el receptor, priorizar la región donde el paciente presenta la curvatura. En pediátricos que no cooperan en PA, realizar en AP.
	
	\textbf{Trucos:} La imagen se invierte horizontalmente al revelar; marcar a la derecha (queda del mismo lado que el corazón). Según criterio Pereira: la marca queda del lado izquierdo y el perfil no se marca. El mayor grado de madurez ósea implica menor posibilidad de corrección sin intervención quirúrgica.
	
	\clearpage
	\onecolumn
	
	% TODO: \usepackage{graphicx} required
	\begin{figure}
		\centering
		\includegraphics[width=0.7\linewidth]{imagenes/col75}
		\caption{}
		\label{fig:col75}
	\end{figure}
	
	
	
	\clearpage
	\twocolumn
	
	%----------------------------------------------------------
	%  Para agregar otro enfoque duplicá el bloque
	%  desde \subsection hasta \textbf{Trucos}
	%---------------------------------------------------------
	
	
%%%%%%%%%%%%%%%%%%%%%%%%%%%%%%%%%%%%%%%%%%%%%%%%%%%%%%%%%%%
%  SECCIÓN: Columna Vertebral — Proyecciones Panorámicas
%  Archivo: secciones/columna_panoramica_lateral.tex
%  Sin preámbulo ni \begin{document} — lo maneja main.tex
	%%%%%%%%%%%%%%%%%%%%%%%%%%%%%%%%%%%%%%%%%%%%%%%%%%%%%%%%%%%
	
	\section{Columna Escoliosis}
	
	\subsection{Panorámica Lateral de Columna (Niño)}
	
	\subsubsection*{1. Identificación}
	
	\textbf{Tipo:} Proyección complementaria. Forma par con la proyección panorámica AP/PA de frente.
	
	\textbf{Indicaciones clínicas principales:}
	\begin{itemize}
		\item Escoliosis: evaluación en el plano sagital.
		\item Espondilolistesis.
		\item Valoración del grado de cifosis y lordosis.
	\end{itemize}
	
	\subsubsection*{2. Técnica Base}
	
	\textbf{Receptor de imagen:} 35x43 cm. Orientación \textbf{longitudinal} (vertical), con respecto al eje mayor de la región. Borde superior a nivel de C7; borde inferior 3--5 cm por debajo de las crestas ilíacas.
	
	\textbf{DFRI:} 180--200 cm (mínimo 150 cm).
	
	\textbf{Técnica:} Homogenizada. Los mAs se duplican respecto al frente para compensar el mayor espesor lateral.
	
	\textbf{Bucky:} Bucky vertical con grilla.
	
	\subsubsection*{3. Posicionamiento}
	
	\paragraph{Paciente}
	De pie (bipedestación), descalzo. Lateralizado al Bucky. El lado \textbf{convexo} de la curva se apoya contra el receptor. Se aplica marcador D o I acorde a hacia donde esté la curvatura.
	
	\paragraph{Región de interés}
	Posición lateral verdadera de pelvis y torso. Brazos llevados hacia adelante a 45°, posición más cómoda para el paciente. Esta posición de brazos evita la superposición sobre la columna y no falsea la lordosis.
	
	\paragraph{Rayo Central}
	Perpendicular al plano del receptor. Punto de entrada: \textbf{4 dedos por dentro del reborde cutáneo posterior}, dirigido al punto medio del receptor de imagen.
	
	\paragraph{Receptor de imagen}
	Borde superior a nivel de C7, dirección caudal. Borde inferior 3--5 cm por debajo de las crestas ilíacas. Debe incluir la región de C7 hacia abajo y al menos 2.5 cm de crestas ilíacas.
	
	\paragraph{Respiración}
	Suspendida al final de la espiración.
	
	\subsubsection*{4. Anatomía Visible}
	
	Vértebras torácicas y lumbares en proyección lateral. Crestas ilíacas (al menos 2.5 cm). En pacientes con escoliosis los arcos costales no estarán bien superpuestos; esto es un hallazgo esperado y no un error técnico, ya que las vértebras de la curvatura pueden estar rotadas, ocasionando distorsión de la caja torácica.
	
	\subsubsection*{5. Criterios de Evaluación}
	
	\textbf{Posicionamiento correcto:} Visualización de C7 hacia abajo con inclusión de al menos 2.5 cm de crestas ilíacas. Posición lateral verdadera de pelvis y torso.
	
	\textbf{Ausencia de rotación:} En lo posible, plano coronal perpendicular al receptor. Se acepta rotación residual inherente a la curvatura escoliótica; los arcos costales no superpuestos en escoliosis dorsal son esperados.
	
	\textbf{Penetración y contraste:} Adecuada visualización de cuerpos vertebrales, espacios intervertebrales y platillos vertebrales en toda la extensión toracolumbar con técnica homogeneizada.
	
	\subsubsection*{6. Puntos Críticos}
	
	\textbf{Elección incorrecta del lado de apoyo:} Siempre realizar primero el frente panorámico para identificar la curvatura. El lado \textbf{cóncavo} debe quedar orientado hacia el tubo: la divergencia cónica del haz hace que los rayos pasen tangenciales a los espacios discales y tiendan a abrirlos. Si se coloca la convexidad hacia el tubo, los espacios articulares se cierran. Ejemplo: concavidad izquierda → perfil derecho (lado izquierdo hacia el tubo).
	
	\textbf{Brazos en posición incorrecta:} Llevar los brazos hacia atrás los superpone a la columna dorsal. Siempre hacia adelante a 45°.
	
	\textbf{Adaptaciones:} En pacientes pediátricos la posición a 45° de brazos resulta más cómoda y reduce el movimiento. Registrar siempre el lado de apoyo utilizado.
	
	\textbf{Trucos:} Consultar siempre la proyección de frente antes de posicionar al paciente para el perfil. La posición de brazos hacia adelante además evita falsear la lordosis lumbar.
	
	
	\onecolumn
	\clearpage
	
	% TODO: \usepackage{graphicx} required
	\begin{figure}
		\centering
		\includegraphics[width=0.3\linewidth]{imagenes/col77}
		\caption{}
		\label{fig:col77}
	\end{figure}
		
	\twocolumn
	\clearpage
	
	%----------------------------------------------------------
	%  Para agregar otro enfoque duplicá el bloque
	%  desde \subsection hasta \textbf{Trucos}
	%----------------------------------------------------------
	
	\subsection{Panorámica Lateral de Columna (Eddy) - Técnica con Doble Exposición}
	
	\subsubsection*{1. Identificación}
	
	\textbf{Tipo:} Proyección complementaria. Forma par con la proyección panorámica AP/PA de frente.
	
	\textbf{Indicaciones clínicas principales:}
	\begin{itemize}
		\item Escoliosis: evaluación en el plano sagital.
		\item Espondilolistesis.
		\item Valoración del grado de cifosis y lordosis.
	\end{itemize}
	
	\subsubsection*{2. Técnica Base}
	
	\textbf{Receptor de imagen:} 35x43 cm. Orientación \textbf{longitudinal} (vertical). Borde superior a nivel de C7; borde inferior 3--5 cm por debajo de las crestas ilíacas.
	
	\textbf{DFRI:} 180--200 cm en ambas exposiciones, para que las dos imágenes queden con igual magnificación y sean comparables.
	
	\textbf{Técnica:} Intermedia. Se realizan \textbf{dos exposiciones} sobre el mismo receptor: la primera centrada en región torácica (C7 hacia abajo, en inspiración) y la segunda centrada en región lumbar (a nivel de crestas ilíacas, en espiración).
	
	\textbf{Bucky:} Bucky vertical con grilla.
	
	\subsubsection*{3. Posicionamiento}
	
	\paragraph{Paciente}
	De pie (bipedestación), descalzo. PMS paralelo al receptor. Lateralizado al Bucky. El lado \textbf{convexo} de la curva primaria se apoya contra el receptor. Se aplica marcador D o I acorde a la curvatura.
	
	% TODO: \usepackage{graphicx} required
	\begin{figure}[h]
		\centering
		\includegraphics[width=0.7\linewidth]{imagenes/col78}
		\caption{}
		\label{fig:col78}
	\end{figure}
	
	% TODO: \usepackage{graphicx} required
	\begin{figure}[h]
		\centering
		\includegraphics[width=0.7\linewidth]{imagenes/col80}
		\caption{}
		\label{fig:col80}
	\end{figure}
	
	
	\paragraph{Región de interés}
	Posición lateral verdadera de pelvis y torso. Brazos llevados hacia arriba, agarrándose por encima de la cabeza, o hacia adelante.
	
	\paragraph{Rayo Central}
	Perpendicular al plano del receptor.
	\begin{itemize}
		\item \textbf{1ra exposición:} punto de entrada 4 dedos por dentro del reborde cutáneo posterior, centrado a nivel de C7 hacia abajo. Para palpar C7, buscar la vértebra prominente. Técnica del ``techito'': mano en la espalda del paciente y por encima la otra mano formando techo con 4 dedos (ajustar si el paciente tiene mayor panículo adiposo).
		\item \textbf{2da exposición:} centrado a la altura de las crestas ilíacas.
	\end{itemize}
	
	\paragraph{Receptor de imagen}
	Borde superior a nivel de C7, dirección caudal. Borde inferior 3--5 cm por debajo de las crestas ilíacas. Debe incluir al menos 2.5 cm de crestas ilíacas.
	
	\paragraph{Respiración}
	\begin{itemize}
		\item \textbf{1ra exposición:} en \textbf{inspiración} (el diafragma baja, libera las vértebras torácicas inferiores).
		\item \textbf{2da exposición:} en \textbf{espiración} (el diafragma sube, libera las vértebras lumbosacras).
	\end{itemize}
	
	\subsubsection*{4. Anatomía Visible}
	
	\textbf{1ra exposición:} vértebras torácicas en proyección lateral.
	\textbf{2da exposición:} vértebras lumbosacras en proyección lateral. En pacientes con escoliosis los arcos costales no estarán bien superpuestos por la rotación vertebral asociada a la curvatura; esto es esperado y no constituye error técnico.
	
	\subsubsection*{5. Criterios de Evaluación}
	
	\textbf{Posicionamiento correcto:} Visualización completa de C7 hacia abajo con inclusión de al menos 2.5 cm de crestas ilíacas. Ambas exposiciones con igual DFRI para comparabilidad.
	
	\textbf{Ausencia de rotación:} PMS paralelo al receptor. Se acepta rotación residual inherente a la escoliosis.
	
	\textbf{Penetración y contraste:} Técnica intermedia con adecuada visualización de cuerpos vertebrales y espacios discales en ambas regiones. Las dos exposiciones deben ser comparables en densidad.
	
	\subsubsection*{6. Puntos Críticos}
	
	\textbf{Diferencia de DFRI entre exposiciones:} Si se varía la distancia entre la primera y segunda exposición, las imágenes quedan con distinta magnificación y no son comparables. Mantener siempre la misma distancia (180 o 200 cm) en ambas.
	
	\textbf{Respiración incorrecta:} Invertir las fases respiratorias oscurece las vértebras de interés. Inspiración para torácicas (diafragma baja), espiración para lumbares (diafragma sube).
	
	\textbf{Adaptaciones:} Si el paciente no puede mantener los brazos por encima de la cabeza, posicionarlos hacia adelante. Ajustar los dedos de referencia del RC si el paciente tiene mayor panículo adiposo.
	
	\textbf{Trucos:} Consultar siempre el frente antes de posicionar: el lado cóncavo de la curvatura debe quedar hacia el tubo para aprovechar la divergencia cónica del haz y abrir los espacios discales. La técnica del
	 ``techito'' con las manos permite localizar el RC con precisión sin necesidad de referencias externas adicionales.
	 
	 \onecolumn
	 \clearpage
	 
	 % TODO: \usepackage{graphicx} required
	 \begin{figure}
	 	\centering
	 	\includegraphics[width=0.3\linewidth]{imagenes/col79}
	 	\caption{Primer disparo}
	 	\label{fig:col79}
	 \end{figure}
	 
	 % TODO: \usepackage{graphicx} required
	 \begin{figure}
	 	\centering
	 	\includegraphics[width=0.3\linewidth]{imagenes/col81}
	 	\caption{Segundo disparo}
	 	\label{fig:col81}
	 \end{figure}
	 
	 
	 \twocolumn
	 \clearpage
	
	%----------------------------------------------------------
	%  Para agregar otro enfoque duplicá el bloque
	%  desde \subsection hasta \textbf{Trucos}
	%----------------------------------------------------------